% See samplepaper.tex and llncsdoc.pdf.
%%% Add the `comments` option to documentclass to enable comments
\documentclass[envcountsame,envcountsect,comments]{llncs}

% conditions for ESOP 2026:
% There is no page limit for submissions, but authors should be aware
% that reviewers are likely to balance the review time for all papers
% and that camera-ready papers may not exceed 25 pages (excluding
% bibliography) and must be formatted in Springer’s LNCS. 

\usepackage[T1]{fontenc} % LNCS.
% \PassOptionsToPackage{square,numbers}{natbib}
% \RequirePackage{natbib}
\usepackage[numbers,compress]{natbib} % LNCS-specific.
\PassOptionsToPackage{inline,shortlabels}{enumitem}
\RequirePackage{enumitem}
\RequirePackage{comment}
\RequirePackage{subcaption}

\PassOptionsToPackage{capitalise,noabbrev}{cleveref}
% LNCS-provided:
% \RequirePackage{amsthm}
% \newtheorem{definition}{Definition}
% \newtheorem{theorem}{Theorem}
% \newtheorem{lemma}{Lemma}

\AtEndPreamble{%
\theoremstyle{acmdefinition}
\newtheorem{remark}[theorem]{Remark}
\newtheorem{convention}[theorem]{Convention}}


\makeatletter
\newcommand*\FreezeMathAlpha[1]{%
  \@tempcnta=#1\relax
  \ifnum\@tempcnta>"7000\relax
    \advance\@tempcnta-"7000\relax
    \mathchardef#1=\@tempcnta
  \fi}
\AtBeginDocument{%
  \FreezeMathAlpha\Gamma  \FreezeMathAlpha\Delta   \FreezeMathAlpha\Theta
  \FreezeMathAlpha\Lambda \FreezeMathAlpha\Xi      \FreezeMathAlpha\Pi
  \FreezeMathAlpha\Sigma  \FreezeMathAlpha\Upsilon \FreezeMathAlpha\Phi
  \FreezeMathAlpha\Psi    \FreezeMathAlpha\Omega}
\makeatother

% https://tex.stackexchange.com/a/40185
\makeatletter
\let\@realLabel\label
\makeatother
\RequirePackage{amsmath}
\RequirePackage{mathtools}

\RequirePackage{xparse}

\RequirePackage{xspace}

% https://tex.stackexchange.com/a/521495
\makeatletter
\@ifdefinable\MakeTokenFromCsname{%
  \long\def\MakeTokenFromCsname#1#{%
    \romannumeral0\@MakeTokenFromCsname{#1}%
  }%
}%
\newcommand\@MakeTokenFromCsname[2]{%
  \expandafter\@@@MakeTokenFromCsname\expandafter{\csname#2\endcsname}{ #1}%
}%
\newcommand\@@@MakeTokenFromCsname[2]{#2#1}%
\makeatother

\makeatletter
\def\DMO#1{%
  \MakeTokenFromCsname\DeclareMathOperator{@#1}{#1}%
  \MakeTokenFromCsname\NewDocumentCommand{#1}{o}{\MakeTokenFromCsname{@#1}\IfValueT{##1}{_{##1}}}%
}
\makeatother

% A little slack for the occasional line that would otherwise stick into
% the margin: applied only to lines TeX cannot break any other way.
\setlength{\emergencystretch}{1em}

\DeclareMathOperator{\fc}{subj}
\DMO{dom}
\DMO{flv}
\DMO{fv}
\DMO{fchans}

\makeatletter

\newcommand{\acronym}[2]{\emph{#1}~(#2)}

\def\hyph-{\mbox{-}\nobreak\hspace{0pt}}
\def\pre#1-{$#1$\mbox{-}\nobreak\hspace{0pt}}
% \nobreak avoids linebreaks
% the 0-width space allows the suffix to be correctly hyphenated

\newcommand{\picalc}{%
    \texorpdfstring{\pre{\pi}-calculus}{pi-calculus}}
\newcommand{\lcalc}{%
    \texorpdfstring{\pre{\lambda}-calculus}{lambda-calculus}}

\makeatother

\newcommand{\Nat}{\mathbb{N}}

\RequirePackage{myproofs}
\RequirePackage{mathpartir}
% Inference rules: mathpartir \infer* inside proofrule(s); see proofrules.sty
\RequirePackage{proofrules}

% `\sepr` for regularly spaced, `\sepr*` for more spacing.
\makeatletter
\def\sepr{\@ifstar\@seprstar\@seprnostar}
\def\@seprnostar{\mathbin{\mbox{\large{$\mid$}}}}
\def\@seprstar{~~\@seprnostar~}
\makeatother

\def\<{\langle}
\def\>{\rangle}
\def\disjoint{\mathrel{\#}}
\NewDocumentCommand\set{m}{{\{#1\}}}
\RequirePackage{stmaryrd}
\NewDocumentCommand\mset{m}{{\lbag#1\rbag}}
\NewDocumentCommand\msetA{m}{{\lbag\begin{array}[t]{l} #1\rbag \end{array}}}
\NewDocumentCommand\update{mmm}{{#1\oplus[#2 \mapsto #3]}}

\let\ltxvdash\vdash
\newcommand{\types}{\mathrel{\color{black}\ltxvdash}}
\newcommand{\proves}{\mathrel{\color{black}\vDash}}

\newcommand{\ichoice}{\mathbin{{\color{cTy}\oplus}}}
\newcommand{\echoice}{\mathrel{{\color{cTy}\&}}}

\PassOptionsToPackage{dvipsnames,rgb}{xcolor}
\RequirePackage{xcolor}

% Semantic markup: sorts, metavariables and marked operators.  Everything
% below marks leaves only; \fFunc/\fType/\fProc survive as aliases.
\RequirePackage{lrlight}
\colorlet{cProc}{cPr}
\colorlet{cType}{cTy}
\colorlet{cFunc}{cTm}
\def\fCtx#1{{\color{cTm}\mathcal{#1}}}

\newcommand{\var}{\mathit}

\def\hole{{\color{cTm}[\cdot]}}

\RequirePackage{upgreek}
% Term keywords carry the term colour; \tyOp is their type-level twin.
\def\tOp#1{\mathinner{{\color{cTm}\mathsf{#1}}}}
\def\tyOp#1{\mathinner{{\color{cTy}\mathsf{#1}}}}
\def\tForce{\tOp{force}}
\def\tFork{\tOp{fork}}
\def\tRecv{\tOp{recv}}
\def\tSend{\tOp{send}}
\def\tWait{\tOp{wait}}
\def\tTerm{\tOp{term}}
\def\tClos{\tOp{clos}}
\def\tLink{\tOp{link}}
\def\tFixOp{\tOp{fix}}
\def\tForce{\tOp{force}}
\def\tInjOp{\tOp{inj}}
\def\tCaseOp{\tOp{case}}
\def\tIfOp{\tOp{if}}
\def\tUnit{{\color{cTm}\ast}}
\def\tyUnit{{\color{cTy}\ast}}
\let\tLVA\vr
\let\tLVB\vrB
\let\tIVX\ivx
\let\tIVY\ivy
\NewDocumentCommand\tAbs{mm}{{{\color{cTm}\uplambda}#1{\color{cTm}.}#2}}
\NewDocumentCommand\tPair{mm}{{\mathopen{{\color{cTm}\<}}#1\mathpunct{{\color{cTm},}}#2\mathclose{{\color{cTm}\>}}}}
\NewDocumentCommand\tBox{m}{{{\color{cTy}\Box}{#1}}}
\NewDocumentCommand\tApp{mm}{{#1\,#2}}
\NewDocumentCommand\tLet{soomm}{{\tOp{let}\IfBooleanTF{#1}{\tUnit}{\tPair{#2}{#3}}=#4\tOp{in}#5}}
\NewDocumentCommand\tFix{mmm}{{\tFixOp\,\tAbs{#1}{\tAbs{#2}{#3}}}}
\NewDocumentCommand\tSub{mmm}{{#1\mathopen{{\color{cTm}[}}#2{\color{cTm}/}#3\mathclose{{\color{cTm}]}}}}
\NewDocumentCommand\tLetI{mmm}{{\tOp{let}{\IExp{#1}}=#2\tOp{in}#3}}
\NewDocumentCommand\tLetS{mmm}{{\tOp{let}{{#1}}=#2\tOp{in}#3}}
\NewDocumentCommand\tFixI{mm}{{\tFixOp\,{#1}.{#2}}}
\NewDocumentCommand\tInj{mm}{{\tInjOp_{#1} #2}}
\NewDocumentCommand\tCase{mmmmm}{{\tCaseOp\,{#1}( #2.#3) (#4.#5)}}
\NewDocumentCommand\tIf{mmm}{{\tIfOp\,{#1}( #2)(#3)}}
\NewDocumentCommand\tSeq{mm}{{#1;#2}}



% Labelled arrows.  \stackon centred the label over a \longrightarrow of
% fixed length, so a long label stuck out at both ends.  \xrightarrow
% grows with its label instead, but bottoms out well short of
% \longrightarrow, so we pad the label to the width that makes the two
% agree: \xrightarrow's width is its label's width plus a fixed overhead,
% which we measure once.
\makeatletter
\newlength\lr@arw@min
\newlength\lr@arw@tmp
\AtBeginDocument{%
  \settowidth\lr@arw@min{$\longrightarrow$}%
  \settowidth\lr@arw@tmp{$\xrightarrow{\mathmakebox[\lr@arw@min]{}}$}%
  \addtolength\lr@arw@tmp{-\lr@arw@min}%  the overhead
  \addtolength\lr@arw@min{-\lr@arw@tmp}%  the label width that matches
  \ifdim\lr@arw@min<\z@ \lr@arw@min\z@ \fi}
\newcommand\lr@arw@lbl[1]{%
  \settowidth\lr@arw@tmp{$\scriptstyle#1$}%
  \ifdim\lr@arw@tmp<\lr@arw@min \lr@arw@tmp\lr@arw@min \fi%
  \mathmakebox[\lr@arw@tmp]{#1}}
\NewDocumentCommand\trans{mo}{%
\ifblank{#1}{\xrightarrow{}}{
  \xrightarrow{\lr@arw@lbl{#1}}
}%
  \IfValueT{#2}{^{#2}}}
% \redu carries no label, but still goes through \xrightarrow so that it
% is the same arrow, and the same length, as \trans with a short one.
\NewDocumentCommand\redu{o}{%
  \longrightarrow\IfValueT{#1}{^{#1}}}
\makeatother
% Step of the label-free reduction semantics (appendix on harmony).
\newcommand{\sred}{\mathrel{\rightsquigarrow}}
\def\irred{\operatorname{irred}}

% Label constructors own their own symbol only; the channels and values
% they mention arrive already marked.
\NewDocumentCommand\lRecv{mm}{{{#1}{?}{#2}}}
\NewDocumentCommand\lSend{mm}{{{#1}{!}{#2}}}
\NewDocumentCommand\lWait{m}{{{#1}{?}}}
\NewDocumentCommand\lTerm{m}{{{#1}{!}}}
\let\lClos\lTerm
\NewDocumentCommand\lFork{mm}{{(\nu{#1}){#2}}}
\NewDocumentCommand\lLink{mm}{{ {\leftrightarrow}\set{{#1}, {#2}}}}
\NewDocumentCommand\lLinkSingle{m}{{ {\leftrightarrow}\set{{#1}}}}

\RequirePackage{amssymb}
\NewDocumentCommand\lFunc{mm}{{{#1}\mathrel{{\color{cTy}\multimap}}{#2}}}
\NewDocumentCommand\lPair{mm}{{{#1}\mathrel{{\color{cTy}\otimes}}{#2}}}
\NewDocumentCommand\lSum{mm}{{{#1}\mathrel{{\color{cTy}+}}{#2}}}
\NewDocumentCommand\sRecv{mm}{{{\color{cTy}?}{#1}{\color{cTy}.}{#2}}}
\NewDocumentCommand\sSend{mm}{{{\color{cTy}!}{#1}{\color{cTy}.}{#2}}}
% \def\sWait{\tOp{Wait}}
% \def\sTerm{\tOp{Term}}
\def\sWait{\tyOp{Close}}
\def\sTerm{\tyOp{Close}}
\def\tEmpty{{\color{cTy}{}\cdot{}}}
\NewDocumentCommand\dual{m}{{\color{cTy}\overline{#1}}}
\def\ptype{\operatorname{ptype}}
\def\primitive{\operatorname{prim}}

\RequirePackage{stmaryrd}
\def\tRepl{\tOp{repl}}
\def\tReq{\tOp{req}}
\makeatletter
% \let\@replSymb\shortuparrow
\let\@replSymb!
\let\@reqSymb\shortdownarrow
\NewDocumentCommand\lRepl{mm}{{{\@replSymb}\lFork{#1}{#2}}}
\NewDocumentCommand\lReq{mm}{{{\@reqSymb}\lFork{#1}{#2}}}
\NewDocumentCommand\lUnres{m}{{{\@replSymb}#1}}
\NewDocumentCommand\IExp{m}{{{\@replSymb}#1}}
\makeatother

\RequirePackage{stmaryrd}
\makeatletter
\NewDocumentCommand\@lr{mmm}{{#1}_{#2}\llbracket{#3}\rrbracket}
\NewDocumentCommand\lrVal{mm}{\@lr{V}{#1}{#2}}
\NewDocumentCommand\lrTerm{mm}{\@lr{E}{#1}{#2}}
\NewDocumentCommand\lrProc{mm}{\@lr{P}{#1}{#2}}
\NewDocumentCommand\lrGSub{mm}{\@lr{G}{#1}{#2}}
% Script T, so the judgment letter does not collide with the
% metavariable T for types.
\NewDocumentCommand\lrTopo{mm}{\@lr{\mathcal{T}}{#1}{#2}}
\makeatother

% Σ ⊩_n (transition) : T  --  "this transition discharges its
% obligation at Σ ⊢ T, with continuations at index n".
\NewDocumentCommand\obl{mmm}{{#1} \mathrel{\Vdash_{#2}} {#3}}

\newcommand{\cupdot}{\mathbin{\mathaccent\cdot\cup}}
% The n-ary form of \cupdot is stmaryrd's \bigcupdot.

\newcommand{\ie}{i.e.,~} % comma
\newcommand{\eg}{e.g.~}  % no comma, according to the Handbook for Scholars

\newcommand{\st}{.\,}

\DeclareMathOperator{\unf}{unfold}
\newcommand{\sflipop}{\mathit{flip}}
\newcommand{\isdual}{\mathrel{\perp}}

\DeclareMathOperator{\sent}{sentc}
% \newcommand{\sent}{sentc}


%%% Local Variables:
%%% mode: latex
%%% TeX-master: "main"
%%% End:


% Load last for the proof environment.
% https://tex.stackexchange.com/a/699868
\let\proof\relax\let\endproof\relax
\usepackage{amsthm}

\RequirePackage{hyperref}
\RequirePackage{cleveref}

% Links from paper claims to the Rocq mechanization.  The table of
% declarations is generated from theories/*.v by paper/rocqrefs.py;
% regenerate with `make -C paper refs`.
\RequirePackage{rocqrefs}
% rocqrefs.sty -- cite the mechanization from the paper, by name.
%
%   \rocq{fundamental_proc}
%       [logo] Compositionality.v:fundamental_proc     (hyperlinked)
%   \rocqas{Theorem 3.13}{fundamental_proc}
%       [logo] Theorem 3.13                            (hyperlinked)
%
% Every reference resolves through paper/rocqrefs.tex, which
% paper/rocqrefs.py regenerates from theories/*.v.  A reference to a
% declaration that no longer exists is a build error, not a dead link.
%
% PERMALINKS.  rocqrefs.py pins the base URL to the HEAD *commit*, so a
% link keeps pointing at the code the paper described even after the
% branch moves on.  Regenerate and the links follow the sources:
%
%     make -C paper refs
%
% ANONYMITY.  The base is the only place the artifact's location
% appears.  For submission regenerate against an Anonymous GitHub
% mirror; nothing else in the paper changes.
%
%     python3 paper/rocqrefs.py gen --base https://anonymous.4open.science/r/XXXX/
%
% LOGO.  Drop the official Rocq logo into paper/ as rocq-logo.pdf (or
% .png) and it is picked up automatically.  Without it the mark falls
% back to a small typeset badge, so the paper builds with no external
% assets and no licensing question.  Override wholesale by redefining
% \rocqlogo.

\NeedsTeXFormat{LaTeX2e}
\ProvidesPackage{rocqrefs}[2026/08/09 links from the paper to the Rocq sources]

\RequirePackage{hyperref}
\RequirePackage{xcolor}
\RequirePackage{graphicx}

\definecolor{rocqlink}{rgb}{0.10,0.35,0.60}
\definecolor{rocqbadge}{rgb}{0.45,0.13,0.13}

% ---- the mark -------------------------------------------------------

% The height is given in ex of the *current* font, so the mark shrinks
% along with the citation text.
\IfFileExists{rocq-logo.pdf}{%
  \newcommand{\rocqlogo}{%
    \raisebox{-0.1ex}{\includegraphics[height=1.15ex]{rocq-logo.pdf}}%
    \kern0.2em}%
}{%
  \IfFileExists{rocq-logo.png}{%
    \newcommand{\rocqlogo}{%
      \raisebox{-0.1ex}{\includegraphics[height=1.15ex]{rocq-logo.png}}%
      \kern0.2em}%
  }{%
    % Fallback badge: a small caps "R" in a rounded rule-free box.
    \newcommand{\rocqlogo}{%
      \raisebox{0.1ex}{\textcolor{rocqbadge}{\textbf{R}}}%
      \kern0.2em}%
  }%
}

% Citations are deliberately set much smaller than the surrounding
% text: they annotate a claim, they are not part of it.
\newcommand{\rocqstyle}{\scriptsize}

% ---- the table ------------------------------------------------------

% Set by the generated file; empty means "typeset names, do not link".
\providecommand{\rocqbase}{}
\newcommand{\rocqsetbase}[1]{\renewcommand{\rocqbase}{#1}}

% \rocqdef{name}{path}{basename}{line} -- one per declaration.
\newcommand{\rocqdef}[4]{%
  \expandafter\gdef\csname rocq@loc@#1\endcsname{#2\#L#4}%
  \expandafter\gdef\csname rocq@file@#1\endcsname{#3}%
}

% Suppress the link border for our links only.  Done with the
% low-level hyperref hook rather than \hypersetup, because these
% citations appear inside captions, where keyval parsing would break.
\newcommand{\rocq@noborder}{%
  \def\@pdfborder{0 0 0}%
  \def\@pdfborderstyle{}%
}

\newcommand{\rocq@fail}[1]{%
  \PackageError{rocqrefs}{No mechanized declaration named `#1'}%
    {Run `make -C paper refs' after changing theories/.}%
}

% ---- the citations --------------------------------------------------

% \rocqas{text}{name}
%
% The link is drawn without a border: the mark and the colored
% monospace text already signal it, and 1000 boxed rules would make
% the figures unreadable.  Scoped to this command, so the document's
% own hyperref setup is untouched.
\newcommand{\rocqas}[2]{%
  \expandafter\ifx\csname rocq@loc@#2\endcsname\relax
    \rocq@fail{#2}%
  \else
    \ifx\rocqbase\empty
      \mbox{\rocqstyle\rocqlogo\texttt{#1}}%
    \else
      \begingroup
        \rocq@noborder
        \href{\rocqbase\csname rocq@loc@#2\endcsname}{%
          \mbox{\rocqstyle\rocqlogo\textcolor{rocqlink}{\texttt{#1}}}}%
      \endgroup
    \fi
  \fi
}

% \rocq{name} -- link text is File.v:name.
\newcommand{\rocq}[1]{%
  \expandafter\ifx\csname rocq@file@#1\endcsname\relax
    \rocq@fail{#1}%
  \else
    \rocqas{\csname rocq@file@#1\endcsname:\detokenize{#1}}{#1}%
  \fi
}

\endinput


% remove in final version:
\hypersetup{
  bookmarksdepth=2,% to get PDF bookmarks back!
  bookmarksopen=true,
  bookmarksopenlevel=1,
}

\usepackage{comments}
% \usepackage{../agda/latex/agda}
\usepackage{txfonts}
% \usepackage{dsfont}
% \usepackage{bbm}
\usepackage{bbold}
\usepackage{newunicodechar}

\DeclareFontEncoding{LS1}{}{}
\DeclareFontSubstitution{LS1}{stix}{m}{n}
\DeclareSymbolFont{symbols2}{LS1}{stixfrak}{m}{n}
\DeclareMathSymbol{\typecolon}{\mathbin}{symbols2}{"25}

\newunicodechar{λ}{\ensuremath{\mathnormal\lambda}}
\newunicodechar{σ}{\ensuremath{\mathnormal\sigma}}
\newunicodechar{τ}{\ensuremath{\mathnormal\tau}}
\newunicodechar{π}{\ensuremath{\mathnormal\pi}}
\newunicodechar{ℕ}{\ensuremath{\mathbb{N}}}
\newunicodechar{∷}{\ensuremath{::}}
\newunicodechar{≡}{\ensuremath{\equiv}}
\newunicodechar{≅}{\ensuremath{\cong}}
\newunicodechar{∀}{\ensuremath{\forall}}
\newunicodechar{ᴸ}{\ensuremath{^L}}
\newunicodechar{ᴿ}{\ensuremath{^R}}
\newunicodechar{ʳ}{\ensuremath{^r}}
\newunicodechar{ⱽ}{\ensuremath{^V}}
\newunicodechar{⟧}{\ensuremath{\rrbracket}}
\newunicodechar{⟦}{\ensuremath{\llbracket}}
\newunicodechar{⊤}{\ensuremath{\top}}
\newunicodechar{⊥}{\ensuremath{\bot}}
\newunicodechar{₁}{\ensuremath{_1}}
\newunicodechar{₂}{\ensuremath{_2}}
\newunicodechar{₃}{\ensuremath{_3}}
\newunicodechar{₄}{\ensuremath{_4}}
\newunicodechar{₅}{\ensuremath{_5}}
\newunicodechar{₆}{\ensuremath{_6}}
\newunicodechar{₇}{\ensuremath{_7}}
\newunicodechar{₈}{\ensuremath{_8}}
\newunicodechar{₉}{\ensuremath{_9}}
\newunicodechar{∈}{\ensuremath{\in}}
\newunicodechar{₀}{\ensuremath{_0}}
\newunicodechar{′}{\ensuremath{'}}
\newunicodechar{ˢ}{\ensuremath{^S}}
\newunicodechar{ᴬ}{\ensuremath{^A}}
\newunicodechar{∘}{\ensuremath{\circ}}
% \newunicodechar{𝟙}{\ensuremath{\mathds{1}}}  
% \newunicodechar{𝟘}{\ensuremath{\mathds{O}}}
% \newunicodechar{𝟙}{\ensuremath{\mathbbm{I}}}  
\newunicodechar{𝟘}{\ensuremath{\mathbb{O}}}
\newunicodechar{ᴾ}{\ensuremath{^P}}
\newunicodechar{ᵀ}{\ensuremath{^T}}
\newunicodechar{⊎}{\ensuremath{\uplus}}
\newunicodechar{ι}{\ensuremath{\iota}}
\newunicodechar{⇐}{\ensuremath{\Leftarrow}}
\newunicodechar{⇒}{\ensuremath{\Rightarrow}}
\newunicodechar{∎}{\ensuremath{\mathnormal\blacksquare}}
\newunicodechar{➙}{\ensuremath{\to^P}}
\newunicodechar{Δ}{\ensuremath{\Delta}}
\newunicodechar{∅}{\ensuremath{\emptyset}}
\newunicodechar{⁺}{\ensuremath{^+}}
\newunicodechar{𝕏}{\ensuremath{\mathbb{X}}}
%\newunicodechar{·}{\ensuremath{\cdot}} %seems to be defined already!
\newunicodechar{∙}{\ensuremath{\bullet}}
\newunicodechar{⁇}{\ensuremath{?}}
\newunicodechar{‼}{\ensuremath{!}}
\newunicodechar{⊕}{\ensuremath{\oplus}}
\newunicodechar{ℤ}{\ensuremath{\mathbb{Z}}}
\newunicodechar{μ}{\ensuremath{\mu}}
\newunicodechar{∃}{\ensuremath{\exists}}
\newunicodechar{⨟}{\ensuremath{\fatsemi}}
\newunicodechar{Σ}{\ensuremath{\Sigma}}
\newunicodechar{ᵣ}{\ensuremath{_r}}
\newunicodechar{ᵢ}{\ensuremath{_i}}
\newunicodechar{≢}{\ensuremath{\nequiv}}
\newunicodechar{≟}{\ensuremath{\stackrel{{\tiny?}}{=}}}
\newunicodechar{≤}{\ensuremath{\le}}
\newunicodechar{ᵇ}{\ensuremath{^b}}
\newunicodechar{𝓣}{\ensuremath{\mathcal{T}}}
\newunicodechar{𝓔}{\ensuremath{\mathcal{E}}}
\newunicodechar{Γ}{\ensuremath{\Gamma}}
\newunicodechar{γ}{\ensuremath{\gamma}}
\newunicodechar{⊔}{\ensuremath{\sqcup}}
\newunicodechar{α}{\ensuremath{\alpha}}
\newunicodechar{η}{\ensuremath{\eta}}
\newunicodechar{ω}{\ensuremath{\omega}}
\newunicodechar{◁}{\ensuremath{\lhd}}
%PT: seems to be already defined
%\newcommand{\lambdabar}{{\mkern0.75mu\mathchar '26\mkern -9.75mu\lambda}}
\newunicodechar{ƛ}{\ensuremath{\lambdabar}}
\newunicodechar{Λ}{\ensuremath{\mathnormal\Lambda}}
\newunicodechar{ρ}{\ensuremath{\rho}}
\newunicodechar{𝓖}{\ensuremath{\mathcal{G}}}
\newunicodechar{ℓ}{\ensuremath{\ell}}
\newunicodechar{♯}{\ensuremath{\sharp}}
\newunicodechar{⇓}{\ensuremath{\Downarrow}}
\newunicodechar{𝓥}{\ensuremath{\mathcal{V}}}
\newunicodechar{∧}{\ensuremath{\wedge}}
\newunicodechar{ₛ}{\ensuremath{_s}}
\newunicodechar{χ}{\ensuremath{\chi}}
\newunicodechar{⊨}{\ensuremath{\models}}
\newunicodechar{⦂}{\ensuremath{\typecolon}}
\newunicodechar{ς}{\ensuremath{\varsigma}}
\newunicodechar{𝓓}{\ensuremath{\mathcal{D}}}
\newunicodechar{∎}{\ensuremath{\square}}
\newunicodechar{■}{\ensuremath{\blacksquare}}
\newunicodechar{↠}{\ensuremath{\twoheadrightarrow}}
\newunicodechar{↪}{\ensuremath{\hookrightarrow}}
\newunicodechar{β}{\ensuremath{\beta}}
\newunicodechar{ξ}{\ensuremath{\xi}}
\newunicodechar{𝟙}{\ensuremath{\mathbb{1}}}
\newunicodechar{𝟚}{\ensuremath{\mathbb{2}}}
\newunicodechar{≔}{\ensuremath{:=}}
\newunicodechar{𝓢}{\ensuremath{\mathcal{S}}}
\newunicodechar{⊢}{\ensuremath{\vdash}}
\newunicodechar{∋}{\ensuremath{\ni}}
\newunicodechar{ₕ}{\ensuremath{_h}}
\newunicodechar{ᵥ}{\ensuremath{_v}}
\newunicodechar{≼}{\ensuremath{\preceq}}
\newunicodechar{𝓗}{\ensuremath{\mathcal{H}}}
\newunicodechar{⊣}{\ensuremath{\dashv}}
\newunicodechar{ᶜ}{\ensuremath{^c}}
\newunicodechar{″}{\ensuremath{''}}
\newunicodechar{‴}{\ensuremath{'''}}
\newunicodechar{ₙ}{\ensuremath{_n}}
\newunicodechar{∣}{\ensuremath{\mid}}
\newunicodechar{⁗}{\ensuremath{''''}}
\newunicodechar{ʰ}{\ensuremath{^h}}
\newunicodechar{✶}{\ensuremath{\ast}}
\newunicodechar{⊸}{\ensuremath{\multimap}}
\newunicodechar{∶}{\ensuremath{\colon}}
\newunicodechar{⊗}{\ensuremath{\otimes}}
\newunicodechar{↭}{\ensuremath{\leftrightsquigarrow}}
\newunicodechar{‵}{\ensuremath{`}}
\newunicodechar{ᴱ}{\ensuremath{^E}}

% \input{../agda-generated}

\title{A Logical-Relations Approach to Compositional Deadlock Freedom for Session Types}
%% (PT: suggest alphabetical order, unsure how to deal with 'van den' and 'd')
%% (BH: Dutch last name sorting ignores 'van den')
%% (ED: Italian last name sorting does *not* ignore the "D'", I am always listed under D. Changed accordingly)
% Emanuele D'Osualdo
% Bas van den Heuvel
% Peter Thiemann 
% Marius Weidner

% \author{%
% Emanuele D'Osualdo\inst{1}\orcidID{0000-0002-9179-5827} \and
% Bas van den Heuvel\inst{2}\orcidID{0000-0002-8264-7371} \and
% Peter Thiemann\inst{3}\orcidID{0000-0002-9000-1239} \and
% Marius Weidner\inst{3}\orcidID{0009-0008-1152-165X}
% }
% %
% % \authorrunning{B. van den Heuvel et al.}

% \institute{%
% University of Konstanz, Germany
% \email{emanuele.dosualdo@uni-konstanz.de}
% \and
% Karlsruhe University of Applied Sciences, Germany
% \and
% University of Freiburg, Germany
% \email{\{thiemann,weidner\}@cs.uni-freiburg.de}
% }

% \NewDocumentCommand\cmnt{mm}{\{\textbf{#1:} #2\}}
% \def\bh{\cmnt{BH}}
% \def\pt{\cmnt{PT}}
% \def\ed{\cmnt{EO}}
% \def\mw{\cmnt{EO}}
\def\OurLang{GV\ensuremath{^{\sharp}}\xspace}

\begin{document}

\maketitle
\pagestyle{plain}

\begin{abstract}
  % Deadlock freedom is a fundamental correctness property of concurrent systems.
  Due to its global nature, deadlock freedom is hard to ensure compositionally.
  Session Types based on a Curry-Howard correspondence with Linear Logic
  achieve compositional deadlock freedom by imposing quite stringent
  constraints on the communication topologies that can be constructed.
  All soundness proofs of such session type systems in the literature
  hinge on topological constraints, making it impossible to safely compose
  typable processes with processes that violate these constraints, but
  which are safe nevertheless.
  
  In this work, we present a new account of session types
  through a logical relation that is
  \begin{enumerate*}[(i)]
    \item Compositional (\ie validates all the type system rules);
    \item Semantic (\ie expresses the meaning of a judgment entirely as constraints on the behavior of untyped processes); and
    \item Liberal (\ie accepts arbitrary processes as long as they are well-behaved on their external interface).
  \end{enumerate*}
  We establish adequacy of the relation for deadlock freedom and
  exercise it on illustrative examples.
  To evaluate the approach, we extend the core calculus with branching, recursion, unrestricted channels and linking.
  Most of our core results are mechanized in Agda.
\end{abstract}

\section{Introduction}
\label{sec:intro}

% PAR: deadlock freedom, global property, compositional typing hard | example: - | figure/rule: - | map: sem-props
% PAR: GV and CP, tree topologies, one channel per cut | example: ex-close | figure/rule: - | map: -
% PAR: priorities, PCP PGV, cyclic structures, numeric annotations | example: ex-cycle | figure/rule: - | map: -
% PAR: our idea, points on ends, order between points, worlds | example: ex-cycle | figure/rule: fig:points-example | map: syn-points, w-world
% PAR: semantic typing, logical relation, typing rules as lemmas, programs outside the type system | example: ex-flag | figure/rule: ex:flag | map: th-compose, th-adeq

\begin{example}[Self-computed flag \rocq{sbv_deadlock_free}]
\label{ex:flag}
% Generated from map example ex-flag (Self-computed flag). Do not edit.
\leavevmode\par\nobreak
\begin{maplisting}
not     = λz. case z of inl w ⇒ inj₂ w | inr w ⇒ inj₁ w
thread  = case not (inj₂ ∗) of
            inl w ⇒ let ∗ = w in let ∗ = close c in close d      -- runs
          | inr w ⇒ let ∗ = w in let ∗ = close d in close c      -- never runs
other   = let ∗ = close  in close 
-- the process ({c, d}, {thread, other}), c at x, d at y
\end{maplisting}

\[
  x \prec \overline{y}\ \text{(running arm)} \qquad
  \overline{x} \prec y\ \text{(other)} \qquad
  y \prec \overline{x}\ \text{(arm that never runs)}
\]
% PAR: running example, thread computes its own flag, E-Case types both arms in one world | example: ex-flag | figure/rule: E-Case, P-Connect | map: r-data, p-typing
% PAR: static rejection, cycle y < xbar < y in the joined world | example: ex-flag | figure/rule: fig:composition-example | map: p-typing, t-lemmas
% PAR: semantic acceptance, thread in E without typing, P-Connect lemma, adequacy | example: ex-flag | figure/rule: thm:compose, thm:adequacy | map: th-compose, th-adeq
\end{example}

% PAR: lock freedom, second level, same figures with coloured additions | example: rx-postpone | figure/rule: fig:typing-core | map: sem-props, th-sound
% PAR: context-free session types in the base system | example: ex-ring | figure/rule: fig:types | map: ty-session

% PAR: contributions | example: - | figure/rule: - | map: -
\begin{itemize}
  \item % PAR: contribution, GV with orders between points, context-free session types, typing rules | example: ex-cycle | figure/rule: sec:typing | map: r-new, r-act, p-typing
  \item % PAR: contribution, compositional logical relation, semantic P-Term and P-Connect, fundamental theorem, adequacy | example: ex-flag | figure/rule: sec:relation | map: th-compose, th-fund-proc, th-adeq
  \item % PAR: contribution, lock freedom by annotations, ranked worlds, fuel | example: rx-postpone | figure/rule: sec:lock-freedom | map: t-lf-counters, t-lf-wf, fu, soup
  \item % PAR: contribution, subsumption of priority and tree systems by example | example: ex-milner | figure/rule: sec:related | map: -
  \item % PAR: contribution, mechanization in Rocq, assumptions EM and FC for lock freedom only | example: - | figure/rule: sec:mechanization | map: th-sound
\end{itemize}

% PAR: outline | example: - | figure/rule: - | map: -


\section{The \OurLang\ Calculus: Syntax, Semantics and Types}
\label{sec:calculus}

\subsection{Syntax and Semantics}

\Cref{fig:syntax} gives the syntax of \OurLang, a concurrent \lcalc\ with
synchronous message passing, loosely based on the formulation of GV by
\citet{conf/esop/LindleyM15}.
Variables are ranged over by $\vr, \vrB$ and channels by $\ch, \chB$.
$\tApp{\tNew}{\tUnit}$ creates a channel and returns both of its ends,
$\tApp{\tSpawn}{\val}$ runs $\tApp{\val}{\tUnit}$ in a new thread,
$\tApp{\tRecv}{\ch}$ receives a message on $\ch$ and returns it with the
continuation channel,
$\tApp{\tSend}{\tPair{\val}{\ch}}$ sends $\val$ on $\ch$,
$\tApp{\tClos}{\ch}$ closes $\ch$,
$\tApp{\tLink}{\tPair{\ch}{\chB}}$ forwards between two ends, and
$\tApp{\tForce}{\tBoxP{\vr}{\tm}}$ unrolls a box.
A box $\tBoxP{\vr}{\tm}$ binds $\vr$ to itself in $\tm$, which is the
only form of recursion in the calculus.

\begin{figure}[thp]
  \adjustfigure[]
  \rulesection*{Constants $\cst$, values $\val$, terms $\tm$ (\rocq{expr})}
  \begin{align*}
    \cst &::=
      \tNew \sepr \tSpawn \sepr \tRecv \sepr \tSend \sepr
      \tClos \sepr \tLink \sepr \tForce
    & \text{(constants)}
    \\
    \val &::=
      \tUnit \sepr \cst \sepr \ch \sepr \tAbs{\vr}{\tm} \sepr
      \tBoxP{\vr}{\tm} \sepr \tPair{\val}{\val} \sepr \tInj{i}{\val}
    & \text{(values)}
    \\
    \tm &::=
      \vr \sepr \val \sepr \tApp{\tm}{\tm} \sepr \tPair{\tm}{\tm} \sepr
      \tLet*{\tm}{\tm} \sepr \tLet[\vr][\vrB]{\tm}{\tm}
    \\ & \qquad {}\sepr
      \tInj{i}{\tm} \sepr \tCase{\tm}{\ivx}{\tm}{\ivy}{\tm}
    & \text{(terms)}
  \end{align*}
  \rulesection{Derived forms}
  \begin{align*}
    \tSeq{\tm_1}{\tm_2} &:= \tLet*{\tm_1}{\tm_2}
    &
    \tAbs{\tUnit}{\tm} &:= \tAbs{\ivx}{\tLet*{\ivx}{\tm}}
    &
    \tLetS{\vr}{\tm_1}{\tm_2} &:= \tApp{(\tAbs{\vr}{\tm_2})}{\tm_1}
  \end{align*}
  \rulesectionend
  \caption{Syntax of \OurLang (\rocq{expr})}
  \Description{Grammars for the constants, values and terms of GV-sharp,
    and three derived forms.}
  \label{fig:syntax}
\end{figure}

\begin{figure}[tp]
  \adjustfigure[]
  \rulesection*{Session labels $\sigma$, term labels $\lambda$, evaluation contexts $\ectx$ (\rocq{slabel}, \rocq{tlabel})}
  \begin{align*}
    \sigma &::=
      \lRecv{\ch}{\val} \sepr
      \lSend{\ch}{\val} \sepr
      \lTerm{\ch} \sepr
      \lLink{\ch}{\chB}
    & \text{(session labels)}
    \\
    \lambda &::=
      \beta \sepr
      \sigma \sepr
      \lNew{\ch} \sepr
      \lSpawnL{\val}
    & \text{(term labels)}
    \\
    \ectx &::=
      \hole \sepr
      \tApp{\ectx}{\tm} \sepr
      \tApp{\val}{\ectx} \sepr
      \tPair{\ectx}{\tm} \sepr
      \tPair{\val}{\ectx} \sepr
      \tLet*{\ectx}{\tm} \sepr
      \tLet[\vr][\vrB]{\ectx}{\tm}
    \\ & \qquad {}\sepr
      \tInj{1}{\ectx} \sepr \tInj{2}{\ectx} \sepr
      \tCase{\ectx}{\ivx_1}{\tm_1}{\ivx_2}{\tm_2}
    & \text{(eval.\ contexts)}
  \end{align*}
  \rulesection{Labelled reductions for terms (\rocq{tstep})}
  \begin{proofrules}
    \infer*[lab=T-Act-App]{}{
      \tApp{(\tAbs{\vr}{\tm})}{\val} \trans{\beta} \tSub{\tm}{\val}{\vr}
    } \label{rule:T-Act-App}
    \and
    \infer*[lab=T-Act-Let*]{}{
      \tLet*{\tUnit}{\tm} \trans{\beta} \tm
    } \label{rule:T-Act-Let*}
    \and
    \infer*[lab=T-Act-LetPair]{}{
      \tLet[\vr][\vrB]{\tPair{\valB}{\val}}{\tm} \trans{\beta} \tSub{\tSub{\tm}{\valB}{\vr}}{\val}{\vrB}
    } \label{rule:T-Act-LetPair}
    \and
    \infer*[lab=T-Act-Case]{}{
      \tCase{(\tInj{\asTm{i}}{\val})}{\ivx_1}{\tm_1}{\ivx_2}{\tm_2} \trans{\beta} \tSub{\tm_{\asTm{i}}}{\val}{\ivx_{\asTm{i}}}
    } \label{rule:T-Act-Case}
    \and
    \infer*[lab=T-Act-Force]{}{
      \tApp{\tForce}{\tBoxP{\vr}{\tm}} \trans{\beta} \tSub{\tm}{\tBoxP{\vr}{\tm}}{\vr}
    } \label{rule:T-Act-Force}
    \and
    \infer*[lab=T-Act-New]{}{
      \tApp{\tNew}{\tUnit} \trans{\lNew{\ch}} \tPair{\ch}{\ch}
    } \label{rule:T-Act-New}
    \and
    \infer*[lab=T-Act-Spawn]{}{
      \tApp{\tSpawn}{\val} \trans{\lSpawnL{\val}} \tUnit
    } \label{rule:T-Act-Spawn}
    \and
    \infer*[lab=T-Act-Recv]{}{
      \tApp{\tRecv}{\ch} \trans{\lRecv{\ch}{\val}} \tPair{\val}{\ch}
    } \label{rule:T-Act-Recv}
    \and
    \infer*[lab=T-Act-Send]{}{
      \tApp{\tSend}{\tPair{\val}{\ch}} \trans{\lSend{\ch}{\val}} \ch
    } \label{rule:T-Act-Send}
    \and
    \infer*[lab=T-Act-Close]{}{
      \tApp{\tClos}{\ch} \trans{\lTerm{\ch}} \tUnit
    } \label{rule:T-Act-Close}
    \and
    \infer*[lab=T-Act-Link]{}{
      \tApp{\tLink}{\tPair{\ch}{\chB}} \trans{\lLink{\ch}{\chB}} \tUnit
    } \label{rule:T-Act-Link}
    \and
    \infer*[lab=T-Act-Ctx]{
      \tm \trans{\lambda} \tm'
    }{
      \ectx[\tm] \trans{\lambda} \ectx[\tm']
    } \label{rule:T-Act-Ctx}
  \end{proofrules}
  \rulesectionend
  \caption{Semantics of \OurLang\ terms}
  \Description{Grammars for evaluation contexts, session labels and term
    labels, and the eleven axioms plus one congruence rule of the
    labelled term reduction relation.}
  \label{fig:term-semantics}
\end{figure}

Both ends of a channel carry one name: \ref{rule:T-Act-New} returns the
pair $\tPair{\ch}{\ch}$. The subjects of a label are
$\fc(\beta) = \fc(\lNew{\ch}) = \fc(\lSpawnL{\val}) = \emptyset$,
$\fc(\lRecv{\ch}{\val}) = \fc(\lSend{\ch}{\val}) = \fc(\lTerm{\ch}) = \set{\ch}$
and $\fc(\lLink{\ch}{\chB}) = \set{\ch,\chB}$, and the channels a label
sends away are $\sent(\lSend{\ch}{\val}) = \fchans(\val)$ and
$\emptyset$ otherwise.
A forwarding label carries an unordered pair, so
$\lLink{\ch}{\chB} = \lLink{\chB}{\ch}$, and
$\tApp{\tLink}{\tPair{\ch}{\ch}}$ emits the singleton
$\lLinkSingle{\ch}$.
We write $\tm \redu \tm'$ for $\tm \trans{\beta} \tm'$ and $\redu[*]$ for
its reflexive transitive closure.


A \emph{process} $\proc$ is a pair $(\chset, \pool)$ of a finite set of
restricted channels $\chset$ and a finite multiset of closed terms
$\pool$ running in parallel.
We write $S_1 \disjoint S_2$ and $S_1 \cupdot S_2$ for disjointness and
disjoint union of sets, and $\pool_1 \uplus \pool_2$ for multiset union.
\Cref{fig:proc-semantics} defines the labelled transition relation
$(\chset,\pool) \trans{\pi} (\chset,\pool)$.
A process may represent a whole network or a sub-network, in which case
the terms of $\pool$ use channels outside $\chset$ to talk to the
environment.
\Cref{rule:P-Act-Session} lifts every offer on a channel shared with the
environment, forwarding included, and removes the sent channels from
$\chset$ in the spirit of the \textsc{Open} rule of the standard LTS for
the \picalc~\cite{lts-pi}.

\begin{remark}
\label{rm:sentc}
Computing the sent channels \emph{syntactically} is not a neutral
choice. In the pure \picalc\ every message is a channel name, while in
\OurLang\ a message is an arbitrary value. Sending
$\val = (\tAbs{\vr}{\asTm{\Omega}\ \ch})$ with $\asTm{\Omega}$ diverging
does not share $\ch$, and still $\ch \in \sent(\val)$, so the LTS
removes it from the local scope. This sound approximation buys a simpler
LTS at the price of two syntactic side conditions in the logical
relation (\cref{rm:vabs}).
\end{remark}

\begin{theorem}[Harmony \rocq{harmony}]
  \label{thm:harmony}
  $\proc \trans{\tau} \proc'$
  if and only if
  $\proc \sred \proc'$
  (\cref{sec:app-harmony}, \cref{fig:red-semantics}).
\end{theorem}

\begin{proof}
  \Cref{sec:app-harmony}, after the reduction semantics it compares the
  LTS with.
\end{proof}

Transitions $\proc \trans{\tau} \proc'$ are written $\proc \redu \proc'$
and called process reductions, with $\trans{}[*]$ for the reflexive
transitive closure.

\begin{figure}[tp]
  \adjustfigure[]
  \rulesection{Labelled reductions for processes (\rocq{plabel})}
  \begin{align*}
    \pi &::= \tau \sepr \sigma
    \qquad \text{(process labels)}
  \end{align*}
  \begin{proofrules}
    \infer*[lab=P-Act-Session]{
      \tm \trans{\sigma} \tm' \\
      \fc(\sigma) \disjoint \chset
    }{
      (\chset, \pool \uplus \mset{\tm}) \trans{\sigma} (\chset \setminus \sent(\sigma), \pool \uplus \mset{\tm'})
    } \label{rule:P-Act-Session}
    \and
    \infer*[lab=P-Act-Beta]{
      \tm \trans{\beta} \tm'
    }{
      (\chset, \pool \uplus \mset{\tm}) \trans{\tau} (\chset, \pool \uplus \mset{\tm'})
    } \label{rule:P-Act-Beta}
    \and
    \infer*[lab=P-Act-New]{
      \ch~\text{fresh}
      \\
      \smash{\tm \trans{\lNew{\ch}} \tm'}
    }{
      (\chset, \pool \uplus \mset{\tm}) \trans{\tau} (\chset \cupdot \set{\ch}, \pool \uplus \mset{\tm'})
    } \label{rule:P-Act-New}
    \and
    \infer*[lab=P-Act-Spawn]{
      \smash{\tm \trans{\lSpawnL{\val}} \tm'}
    }{
      (\chset, \pool \uplus \mset{\tm}) \trans{\tau} (\chset, \pool \uplus \mset{\tm', \tApp{\val}{\tUnit}})
    } \label{rule:P-Act-Spawn}
    \and
    \infer*[lab=P-Act-Msg]{
      \tm \trans{\lSend{\ch}{\val}} \tm' \\
      \tmB \trans{\lRecv{\ch}{\val}} \tmB'
    }{
      (\chset, \pool \uplus \mset{\tm, \tmB}) \trans{\tau} (\chset, \pool \uplus \mset{\tm', \tmB'})
    } \label{rule:P-Act-Msg}
    \and
    \infer*[lab=P-Act-Close]{
      \tm \trans{\lTerm{\ch}} \tm' \\
      \tmB \trans{\lTerm{\ch}} \tmB' \\
      \ch \in \chset
    }{
      (\chset, \pool \uplus \mset{\tm, \tmB}) \trans{\tau} (\chset \setminus \set{\ch}, \pool \uplus \mset{\tm', \tmB'})
    } \label{rule:P-Act-Close}
    \and
    \infer*[lab=P-Act-Link]{
      \tm \trans{\lLink{\ch}{\chB}} \tm' \\
      \ch \in \chset
    }{
      (\chset, \pool \uplus \mset{\tm}) \trans{\tau} (\chset \setminus \set{\ch}, \tSub{(\pool \uplus \mset{\tm'})}{\chB}{\ch})
    } \label{rule:P-Act-Link}
  \end{proofrules}
  \rulesectionend
  \caption{LTS semantics of \OurLang\ processes}
  \Description{The grammar of process labels and the seven rules of the
    process labelled transition relation: one lifting an external offer,
    and one each for beta steps, channel creation, spawning, matched
    messages, matched closes and forwarding.}
  \label{fig:proc-semantics}
\end{figure}

\paragraph{Properties of processes}
Forked threads are never joined, so they are expected to reduce to unit,
while the initial thread may evaluate to the value computed by the whole
program.\footnote{This design decision is from
  \citet{conf/esop/LindleyM15}. The rest of the development does
  \emph{not} depend on it.}

\begin{definition}[Irreducible terms and processes \rocq{irred}]
We write $\irred \tm$ if there are no $\lambda, \tm'$ with
$\tm \trans{\lambda} \tm'$, and $\irred \proc$ if there are no
$\pi, \proc'$ with $\proc \trans{\pi} \proc'$.
\end{definition}

\begin{definition}[Terminated process \rocq{terminated}]
\label{def:terminated}
  A process $\proc$ is \emph{terminated} if
  $\proc = (\asPrOp{\emptyset}, \mset{\val,\tUnit,\dots,\tUnit})$
  for some~$\val$.
\end{definition}

\begin{definition}[Safe process \rocq{safe}]
\label{def:safe}
  A process $\proc$ is \emph{safe} if it is terminated, or if
  $\proc \redu \proc'$ for some~$\proc'$.
\end{definition}

\begin{definition}[Deadlock freedom \rocq{deadlock_free}]
  \label{def:df}
  A process $\proc$ is \emph{deadlock free} if
  ${\forall \proc'.\,\proc \redu[*] \proc'}$ implies $\proc'$ is safe.
\end{definition}

\begin{example}[An exchange between two threads \rocq{sync_main}]
\label{ex:sync}
  \begin{align*}
    \asTmId{sync} := {}
      & \tLetS{\tPair{\ch}{\chB}}{\tApp{\tNew}{\tUnit}}{}
        \tSeq{\tApp{\tSpawn}{(\tAbs{\tUnit}{\tApp{\tClos}{(\tApp{\tSend}{\tPair{\tUnit}{\ch}})}})}}{}
      \\
      & \tLet[\ivx][\chB]{\tApp{\tRecv}{\chB}}{\tSeq{\ivx}{\tApp{\tClos}{\chB}}}
  \end{align*}
  The spawned thread sends $\tUnit$ on $\ch$ and closes it, the main
  thread receives on $\chB$ and closes it. The process
  $(\emptyset, \mset{\asTmId{sync}})$ has exactly one reduction sequence,
  and it ends terminated, so it is safe and deadlock free
  (\rocq{sync_deadlock_free}).
\end{example}

\begin{example}[A cyclic topology \rocq{cycle_main}]
\label{ex:cyclic}
  \begin{align*}
    \asTmId{cycle} := {}
      & \tLetS{\tPair{\ch}{\ch'}}{\tApp{\tNew}{\tUnit}}{}
        \tLetS{\tPair{\chB}{\chB'}}{\tApp{\tNew}{\tUnit}}{}
      \\
      & \tSeq{\tApp{\tSpawn}{(\tAbs{\tUnit}{\tSeq{\tApp{\tClos}{\ch'}}{\tApp{\tClos}{\chB'}}})}}{}
        \tSeq{\tApp{\tClos}{\ch}}{\tApp{\tClos}{\chB}}
  \end{align*}
  Two threads share the two channels $\ch$ and $\chB$, so the topology
  is cyclic. Both threads close $\ch$ before $\chB$, the process has
  exactly one reduction sequence, and it is deadlock free
  (\rocq{cycle_deadlock_free}).
\end{example}

\begin{example}[A deadlocked process \rocq{blocked_main}]
\label{ex:deadlock}
  Reversing the two closes of one thread of \cref{ex:cyclic} gives
  \begin{align*}
    \asTmId{blocked} := {}
      & \tLetS{\tPair{\ch}{\ch'}}{\tApp{\tNew}{\tUnit}}{}
        \tLetS{\tPair{\chB}{\chB'}}{\tApp{\tNew}{\tUnit}}{}
      \\
      & \tSeq{\tApp{\tSpawn}{(\tAbs{\tUnit}{\tSeq{\tApp{\tClos}{\chB'}}{\tApp{\tClos}{\ch'}}})}}{}
        \tSeq{\tApp{\tClos}{\ch}}{\tApp{\tClos}{\chB}}
  \end{align*}
  After the two creations and the spawn the process is irreducible and
  not terminated: one thread waits to close $\ch$, the other waits to
  close $\chB$. It is not safe, so no process that reaches it is
  deadlock free.
\end{example}

\subsection{The Static System}
\label{sec:syntacticTyping}
\label{sec:syntacticTyping:typeSys}

Types carry \emph{points}. A point is a name with a natural offset, a
channel type is $\sAt{\sty}{\pntA}$, and the $k$-th action of a channel
at $\pntA$ happens at $\poff{\pntA}{k}$. \Cref{fig:points} gives points,
facts and constraints, \cref{fig:types} the types, and
\cref{fig:facts} the fact theory the rules read every side condition
from.

\begin{figure}[tp]
  \adjustfigure[]
  \rulesection*{Names $\pnm$, points $\pntA$, facts $\fact$, constraints $\cstr$ (\rocq{point}, \rocq{fact}, \rocq{bcs})}
  \begin{align*}
    \pnm &::= \pcst \sepr \pviA \sepr \pslot{\pntA}{j}
    & \text{(a constant, a point variable, the $j$-th channel carried at $\pntA$)}
    \\
    \pntA &::= \poff{\pnm}{k}
    & \text{(the $k$-th action of $\pnm$, $k \in \Nat$; $\poff{\pnm}{0}$ is written $\pnm$)}
    \\
    \fact \in \tyFact &::= \fLt{\pntA}{\pntB} \sepr \fEq{\pntA}{\pntB}
    & \text{($\pntA$ acts before $\pntB$; $\pntA$ and $\pntB$ are one point)}
    \\
    \cstr &\subseteq \tyFact
    & \text{(a set of facts, finite or infinite, under a binder $\pvivec$)}
  \end{align*}
  \rulesection{Roots}
  \begin{align*}
    \rootof(\poff{\pviA}{k}) &= \pviA
    &
    \rootof(\poff{\pslot{\pntA}{j}}{k}) &= \rootof(\pntA)
    &
    \rootof(\poff{\pcst}{k}) &= \pcst
  \end{align*}
  \rulesectionend
  \caption{Points, facts and constraints}
  \Description{The grammars of names, points, facts and constraints, and
    the three equations defining the root of a point.}
  \label{fig:points}
\end{figure}

A channel at $\pntA$ acts there and moves to $\poff{\pntA}{1}$. A
message sent at $\pntA$ carries its $j$-th channel at $\pslot{\pntA}{j}$.
Every point of a judgment is built from a point variable or a constant.
$\pvivec$ stands for a sequence $\pviA_1 \dots \pviA_m$ of point
variables and $\pntvec$ for a sequence of points, and
$\tysub{\ty}{\pntvec}{\pvivec}$ replaces each $\pviA_i$ by $\pntA_i$.

\begin{figure}[tp]
  \adjustfigure[]
  \rulesection*{Types $\ty$, session types $\sty$, bounds $\ubd, \lbd$ (\rocq{ty})}
  \begin{align*}
    \ty,\tyB &::= \tyUnit \sepr \lPair{\ty}{\tyB} \sepr \lSum{\ty}{\tyB}
      \sepr \ty \lolli{\ubd}{\lbd} \tyB \sepr \sAt{\sty}{\pntA}
      \sepr \tyBox{\pvivec}{\cstr}{\ty} \sepr \tyEx{\lbd}{\pvivec}{\cstr}{\ty}
    \\
    \sty &::= \sSend{\ty}{\sty} \sepr \sRecv{\ty}{\sty} \sepr \sTerm
      \sepr \tyMu\tvar.\sty \sepr \tvar \sepr \dualv{\tvar}
    \\
    \ubd &::= \pBot \sepr \pntA
    & \lbd &::= \pTop \sepr \pBot \sepr \pntA
  \end{align*}
  \rulesection{Session duality $\dual{\sty}$ (\rocq{dual})}
  \begin{align*}
    \dual{\sSend{\ty}{\sty}} &= \sRecv{\ty}{\dual{\sty}}
    &
    \dual{\sRecv{\ty}{\sty}} &= \sSend{\ty}{\dual{\sty}}
    &
    \dual{\sTerm} &= \sTerm
    \\
    \dual{\tvar} &= \dualv{\tvar}
    &
    \dual{\dualv{\tvar}} &= \tvar
    &
    \dual{\tyMu\tvar.\sty} &= \tyMu\tvar.\,\sflipop_{\tvar}(\dual{\sty})
  \end{align*}
  \rulesection{Point substitution $\tysub{\pntA}{\pntvec}{\pvivec}$, $\tysub{\ty}{\pntvec}{\pvivec}$ (\rocq{pt_inst}, \rocq{pinst_ty})}
  \begin{align*}
    \tysub{(\poff{\pviA_i}{k})}{\pntvec}{\pvivec} &= \poff{\pntA_i}{k}
    &
    \tysub{(\poff{\pcst}{k})}{\pntvec}{\pvivec} &= \poff{\pcst}{k}
    &
    \tysub{(\poff{\pslot{\pntB}{j}}{k})}{\pntvec}{\pvivec}
      &= \poff{\pslot{(\tysub{\pntB}{\pntvec}{\pvivec})}{j}}{k}
    \\
    \tysub{(\sAt{\sty}{\pntB})}{\pntvec}{\pvivec}
      &= \sAt{\sty}{(\tysub{\pntB}{\pntvec}{\pvivec})}
    \span\span\span\span
  \end{align*}
  and structurally on the other type formers, mapping over the bounds
  $\ubd$ and $\lbd$ of an arrow and of an existential, and going under
  the binder of a $\Box$ and of an $\exists$ (\rocq{pup}).
  \rulesection{Payloads (\rocq{payload_at})}
  \begin{align*}
    \tysub{\ty}{\pntA}{\pself} &= \tysub{\ty}{\pslot{\pntA}{j}}{\pslot{\pself}{j}}
  \end{align*}
  \rulesectionend
  \caption{Types, session types and duality}
  \Description{The grammars of types, session types and the two kinds of
    bound; the six equations of session duality; and the definition of a
    payload instantiated at a point.}
  \label{fig:types}
\end{figure}

A channel of type $\sAt{\sty}{\pntA}$ runs $\sty$ and acts next at
$\pntA$. The arrow $\ty \lolli{\ubd}{\lbd} \tyB$ carries the latest
point at which a call may block, $\pBot$ for never, and a lower bound on
the points of the channels the closure holds, $\pTop$ for none. Arrows
take $\lbd \in \set{\pTop} \cup \{\pntA\}$ and only existentials take
$\lbd = \pBot$. $\tyBox{\pvivec}{\cstr}{\ty}$ is replicable, holds no
channel, and is usable at every choice of points for $\pvivec$ that
satisfies $\cstr$. $\tyEx{\lbd}{\pvivec}{\cstr}{\ty}$ is $\ty$ at some
such choice.

\begin{definition}[Payload \rocq{payload_at}]
\label{def:payload}
  The type $\ty$ in $\sSend{\ty}{\sty}$ and $\sRecv{\ty}{\sty}$ mentions
  no point other than the markers $\pslot{\pself}{0}, \pslot{\pself}{1},
  \dots$, and a send or receive at $\pntA$ fills $\pslot{\pself}{j}$
  with $\pslot{\pntA}{j}$, written $\tysub{\ty}{\pntA}{\pself}$.
\end{definition}

\begin{remark}
\label{rm:payload}
  A marker of a payload is a self-marker by construction, so no rule
  enforces \cref{def:payload}. Point substitution replaces the
  variables $\pvivec$, and $\pslot{\pself}{j}$ is no variable, so no
  enclosing $\Box$ or $\exists$ touches it. Point substitution also
  stops at a channel type, by the clause
  $\tysub{(\sAt{\sty}{\pntB})}{\pntvec}{\pvivec} =
    \sAt{\sty}{(\tysub{\pntB}{\pntvec}{\pvivec})}$ of \cref{fig:types},
  so it never reaches inside a payload at all (\rocq{pinst_ty}).

  Session substitution does reach inside a payload, because unfolding
  $\tyMu\tvar.\sRecv{(\sAt{\tvar}{\pslot{\pself}{0}})}{\sTerm}$ has to
  rewrite the $\tvar$ under the carrier. It rewrites session variables
  and nothing else, and the point a channel type hangs on stays put
  (\rocq{sinst_ty}, and the same shape in \rocq{sren_ty} and
  \rocq{sflip_ty}), so unfolding leaves every marker alone. The
  property the proofs use is
  $\tysub{(\tysub{\ty}{\pntA}{\pself})}{\pntvec}{\pvivec}
    = \tysub{\ty}{(\tysub{\pntA}{\pntvec}{\pvivec})}{\pself}$
  (\rocq{payload_pinst}), which holds with no side condition because
  $\tysub{\ty}{\pntA}{\pself}$ is a total substitution.
\end{remark}

Sender and receiver name the carried channels alike, so a message
carries no point. $\dualv{\tvar}$ is the dual polarity of $\tvar$, the
dual of whatever $\tvar$ stands for. Source types never contain
$\dualv{\tvar}$: it arises only from duality, and dualizing
$\tyMu\tvar.\sty$ swaps the polarities of its own variable so that
unfolding commutes with duality (\cref{sec:app-duality}).

\paragraph{Contexts and the judgment}
A linear environment $\lenv ::= \tEmpty \sepr \lenv, \vr : \ty$ is
linear except at $\Box$ types. $\lenv[\vr : \ty]$ updates $\lenv$, which
already binds $\vr$, so that $\vr$ has type $\ty$, and it extends
nothing. The facts $\fctx$ assumed by a judgment
are facts $\fact$ over points together with facts $\fLt{\ubd}{\pfv}$ for
frame variables $\pfv$. Nodes are
$\pnode ::= \pntA \sepr \pfv \sepr \pBot$, and a frame $\frm$ is a
finite set of nodes. The judgment
\[
  \fctx; \lenv \typesF{\frm} \tm : \ty
  \qquad(\rocq{styped})
\]
says that $\tm$ uses exactly $\lenv$ and has type $\ty$, and that the
rest of the program holds the nodes of $\frm$ while $\tm$ runs. A frame
variable $\pfv$ stands for whatever the caller of a function holds, and
$\pBot$ for channels at points not known here.

\begin{figure}[tp]
  \adjustfigure[]
  \rulesection*{Aliasing $\fctx \types \fEq{\pntA}{\pntB}$ (\rocq{seqs})}
  \begin{proofrules}
    \infer*[lab=A-Hyp]{
      (\fEq{\pntA}{\pntB}) \in \fctx
    }{
      \fctx \types \fEq{\pntA}{\pntB}
    } \label{rule:A-Hyp}
    \and
    \infer*[lab=A-Refl]{}{
      \fctx \types \fEq{\pntA}{\pntA}
    } \label{rule:A-Refl}
    \and
    \infer*[lab=A-Sym]{
      \fctx \types \fEq{\pntA}{\pntB}
    }{
      \fctx \types \fEq{\pntB}{\pntA}
    } \label{rule:A-Sym}
    \and
    \infer*[lab=A-Trans]{
      \fctx \types \fEq{\pntA}{\pntB} \\
      \fctx \types \fEq{\pntB}{\pntC}
    }{
      \fctx \types \fEq{\pntA}{\pntC}
    } \label{rule:A-Trans}
    \and
    \infer*[lab=A-Shift]{
      \fctx \types \fEq{\pntA}{\pntB}
    }{
      \fctx \types \fEq{\poff{\pntA}{1}}{\poff{\pntB}{1}}
    } \label{rule:A-Shift}
    \and
    \infer*[lab=A-Slot]{
      \fctx \types \fEq{\pntA}{\pntB}
    }{
      \fctx \types \fEq{\pslot{\pntA}{j}}{\pslot{\pntB}{j}}
    } \label{rule:A-Slot}
  \end{proofrules}
  \rulesection{Precedence $\fctx \types \fLt{\pnode}{\pnode'}$ (\rocq{fl})}
  \begin{proofrules}
    \infer*[lab=O-Po]{}{
      \fctx \types \fLt{\pntA}{\poff{\pntA}{1}}
    } \label{rule:O-Po}
    \and
    \infer*[lab=O-Hyp]{
      (\fLt{\pnode}{\pnode'}) \in \fctx
    }{
      \fctx \types \fLt{\pnode}{\pnode'}
    } \label{rule:O-Hyp}
    \and
    \infer*[lab=O-EqL]{
      \fctx \types \fEq{\pntA}{\pntA'} \\
      \fctx \types \fLt{\pntA'}{\pnode}
    }{
      \fctx \types \fLt{\pntA}{\pnode}
    } \label{rule:O-EqL}
    \and
    \infer*[lab=O-EqR]{
      \fctx \types \fLt{\pnode}{\pntB'} \\
      \fctx \types \fEq{\pntB'}{\pntB}
    }{
      \fctx \types \fLt{\pnode}{\pntB}
    } \label{rule:O-EqR}
    \and
    \infer*[lab=O-Trans]{
      \fctx \types \fLt{\pnode}{\pntB} \\
      \fctx \types \fLt{\pntB}{\pnode'}
    }{
      \fctx \types \fLt{\pnode}{\pnode'}
    } \label{rule:O-Trans}
  \end{proofrules}
  \rulesection{Heads, lower bounds, precedence of a frame (\rocq{heads})}
  \begin{align*}
    \heads(\sAt{\sty}{\pntA}) &= \set{\pntA}
    &
    \heads(\ty \lolli{\ubd}{\lbd} \tyB) &= \heads(\tyEx{\lbd}{\pvivec}{\cstr}{\ty}) = \set{\lbd}
    \\
    \heads(\lPair{\ty}{\tyB}) &= \heads(\lSum{\ty}{\tyB}) = \heads(\ty) \cup \heads(\tyB)
    \\
    \heads(\tyUnit) &= \heads(\tyBox{\pvivec}{\cstr}{\ty}) = \emptyset
  \end{align*}
  with $\pTop$ contributing nothing, and
  $\heads(\lenv) = \bigcup \set{\heads(\ty) \mid \vr : \ty \in \lenv}$.
  \begin{align*}
    \fctx \types \fLe{\lbd}{N}
    &\quad\text{iff}\quad
    \lbd = \pTop \text{ and } N = \emptyset,
    \text{ or } \lbd = \pBot,
    \\ & \qquad\qquad
    \text{ or } \lbd = \pntA \text{ and }
      \fctx \types \fEq{\pntA}{\pnode} \text{ or } \fctx \types \fLt{\pntA}{\pnode}
      \text{ for every } \pnode \in N
    \\
    \fctx \types \fLt{\ubd}{\frm}
    &\quad\text{iff}\quad
    \ubd = \pBot,
    \text{ or } \ubd = \pntA \text{ and }
      \fctx \types \fLt{\pntA}{\pnode} \text{ for every } \pnode \in \frm
  \end{align*}
  \rulesectionend
  \caption{The fact theory}
  \Description{Six rules for aliasing, five rules for precedence, the
    heads of a type and of an environment, and the two derived
    judgments relating a bound to a set of nodes.}
  \label{fig:facts}
\end{figure}

\Cref{rule:O-Po} is program order, and it relates two points of one name
only. No rule derives $\fLt{\pnode}{\pBot}$, so a $\pBot$ node in $N$ or
in $\frm$ makes both derived judgments fail for a point: no closure,
message, packed value or blocking call coexists with an unopened
existential at $\pBot$.

\paragraph{Term rules}
\Cref{fig:typeSys,fig:typeSysPoints} give the term rules. $\lenv_1 \cdot \lenv_2$ splits
$\lenv$ and copies variables of $\Box$ type to both sides. The first
subterm of every composite runs while the rest holds its channels, so it
is typed under $\frm$ grown by the heads of that rest.
\ref{rule:E-Abs} types the body under the frame $\set{\pfv}$ and assumes
$\fLt{\ubd}{\pfv}$, so every call in the body precedes $\ubd$, and
\ref{rule:E-App} checks that $\ubd$ precedes what the caller holds. The
two checks meet at $\ubd$.
\ref{rule:E-Exp} is the one recursive rule and the only rule that binds
$\forall$: $\forall$ appears only directly under $\Box$, because a value
that holds channels has fixed points.

\begin{figure}[tp]
  \adjustfigure[]
  \rulesection*{Typing rules for terms $\fctx ; \lenv \typesF{\frm} \tm : \ty$ (\rocq{styped})}
  \begin{proofrules}
    \infer*[lab=E-Unit]{}{
      \fctx ; \tEmpty \typesF{\frm} \tUnit : \tyUnit
    } \label{rule:E-Unit}
    \and
    \infer*[lab=E-Const]{
      \cst : \ty \text{ an instance of a signature under } \fctx
    }{
      \fctx ; \tEmpty \typesF{\frm} \cst : \ty
    } \label{rule:E-Const}
    \and
    \infer*[lab=E-Var]{}{
      \fctx ; \vr : \ty \typesF{\frm} \vr : \ty
    } \label{rule:E-Var}
    \and
    \infer*[lab=E-Abs]{
      \pfv \text{ fresh} \\
      \fctx \types \fLe{\lbd}{\heads(\lenv)} \\
      \fctx, \fLt{\ubd}{\pfv} ; \lenv , \vr : \ty \typesF{\set{\pfv}} \tm : \tyB
    }{
      \fctx ; \lenv \typesF{\frm} \tAbs{\vr}{\tm} : \ty \lolli{\ubd}{\lbd} \tyB
    } \label{rule:E-Abs}
    \and
    \infer*[lab=E-App]{
      \fctx \types \fLt{\ubd}{\frm} \\
      \fctx ; \lenv_1 \typesF{\frm \cup \heads(\lenv_2)} \tm_1 : \ty \lolli{\ubd}{\lbd} \tyB \\
      \fctx ; \lenv_2 \typesF{\frm \cup \set{\lbd}} \tm_2 : \ty
    }{
      \fctx ; \lenv_1 \cdot \lenv_2 \typesF{\frm} \tApp{\tm_1}{\tm_2} : \tyB
    } \label{rule:E-App}
    \and
    \infer*[lab=E-Pair]{
      \fctx ; \lenv_1 \typesF{\frm \cup \heads(\lenv_2)} \tm_1 : \ty \\
      \fctx ; \lenv_2 \typesF{\frm \cup \heads(\ty)} \tm_2 : \tyB
    }{
      \fctx ; \lenv_1 \cdot \lenv_2 \typesF{\frm} \tPair{\tm_1}{\tm_2} : \lPair{\ty}{\tyB}
    } \label{rule:E-Pair}
    \and
    \infer*[lab=E-LetPair]{
      \fctx ; \lenv_1 \typesF{\frm \cup \heads(\lenv_2)} \tm_1 : \lPair{\tyB}{\tyC} \\
      \fctx ; \lenv_2 , \vr : \tyB , \vrB : \tyC \typesF{\frm} \tm_2 : \ty
    }{
      \fctx ; \lenv_1 \cdot \lenv_2 \typesF{\frm} \tLet[\vr][\vrB]{\tm_1}{\tm_2} : \ty
    } \label{rule:E-LetPair}
    \and
    \infer*[lab=E-Let*]{
      \fctx ; \lenv_1 \typesF{\frm \cup \heads(\lenv_2)} \tm_1 : \tyUnit \\
      \fctx ; \lenv_2 \typesF{\frm} \tm_2 : \ty
    }{
      \fctx ; \lenv_1 \cdot \lenv_2 \typesF{\frm} \tLet*{\tm_1}{\tm_2} : \ty
    } \label{rule:E-Let*}
    \and
    \infer*[lab=E-Inj]{
      \fctx ; \lenv \typesF{\frm} \tm : \ty_{\asTy{i}}
    }{
      \fctx ; \lenv \typesF{\frm} \tInj{\asTm{i}}{\tm} : \lSum{\ty_1}{\ty_2}
    } \label{rule:E-Inj}
    \and
    \infer*[lab=E-Case]{
      \fctx ; \lenv_1 \typesF{\frm \cup \heads(\lenv_2)} \tm : \lSum{\ty_1}{\ty_2} \\
      \fctx ; \lenv_2 , \ivx_1 : \ty_1 \typesF{\frm} \tm_1 : \tyB \\
      \fctx ; \lenv_2 , \ivx_2 : \ty_2 \typesF{\frm} \tm_2 : \tyB
    }{
      \fctx ; \lenv_1 \cdot \lenv_2 \typesF{\frm} \tCase{\tm}{\ivx_1}{\tm_1}{\ivx_2}{\tm_2} : \tyB
    } \label{rule:E-Case}
  \end{proofrules}
  \rulesectionend
  \caption{Typing rules for terms: variables, functions and data}
  \Description{Ten term typing rules, for the unit, constants,
    variables, abstraction, application, pairs, their elimination,
    sequencing, injections and case.}
  \label{fig:typeSys}
\end{figure}

\begin{figure}[tp]
  \adjustfigure[]
  \rulesection*{Structural rules and rules for points, $\fctx ; \lenv \typesF{\frm} \tm : \ty$ (\rocq{styped})}
  \begin{proofrules}
    \infer*[lab=E-Weak]{
      \fctx ; \lenv \typesF{\frm} \tm : \ty \\
      \lenv' \text{ only } \Box \text{ types}
    }{
      \fctx ; \lenv , \lenv' \typesF{\frm} \tm : \ty
    } \label{rule:E-Weak}
    \and
    \infer*[lab=E-Contract]{
      \fctx ; \lenv , \ivx : \tBox{\ty} , \ivy : \tBox{\ty} \typesF{\frm} \tm : \tyB
    }{
      \fctx ; \lenv , \ivx : \tBox{\ty} \typesF{\frm} \tSub{\tm}{\ivx}{\ivy} : \tyB
    } \label{rule:E-Contract}
    \and
    \infer*[lab=E-Conv]{
      \fctx ; \lenv \typesF{\frm} \tm : \ty \\
      \ty \approx \ty' \text{ up to } \fctx \types {\equiv}
    }{
      \fctx ; \lenv \typesF{\frm} \tm : \ty'
    } \label{rule:E-Conv}
    \and
    \infer*[lab=E-Exp]{
      \lenv \text{ only } \Box \text{ types} \\
      \pvivec, \pfv \text{ fresh} \\\\
      \fctx, \cstr, \fLt{\ubd}{\pfv} ;
        \lenv , \vr : \tyBox{\pvivec}{\cstr}{\tyUnit \lolli{\ubd}{\pTop} \ty}
        \typesF{\set{\pfv}} \tm : \ty
    }{
      \fctx ; \lenv \typesF{\frm} \tBoxP{\vr}{\tm}
        : \tyBox{\pvivec}{\cstr}{\tyUnit \lolli{\ubd}{\pTop} \ty}
    } \label{rule:E-Exp}
    \and
    \infer*[lab=E-Pack]{
      \fctx \types \tysub{\cstr}{\pntvec}{\pvivec} \\
      \fctx \types \fLe{\lbd}{\heads(\tysub{\ty}{\pntvec}{\pvivec})} \\
      \fctx ; \lenv \typesF{\frm} \tm : \tysub{\ty}{\pntvec}{\pvivec}
    }{
      \fctx ; \lenv \typesF{\frm} \tm : \tyEx{\lbd}{\pvivec}{\cstr}{\ty}
    } \label{rule:E-Pack}
    \and
    \infer*[lab=E-Open]{
      \vr : \tyEx{\lbd}{\pvivec}{\cstr}{\ty} \in \lenv \\
      \pvivec \text{ fresh, not in } \tyB \\
      \fctx, \cstr ; \lenv[\vr : \ty] \typesF{\frm} \tm : \tyB
    }{
      \fctx ; \lenv \typesF{\frm} \tm : \tyB
    } \label{rule:E-Open}
  \end{proofrules}
  \rulesectionend
  \caption{Typing rules for terms: structure and points}
  \Description{Six term typing rules, for weakening, contraction,
    conversion, a box, packing and opening an existential.}
  \label{fig:typeSysPoints}
\end{figure}

\paragraph{Signatures}
\Cref{fig:signatures} gives the type-indexed family of signatures the
constants are instantiated from. An action that may block happens at a
point $\pntA$, so $\tSend$, $\tRecv$ and $\tClos$ have upper bound
$\pntA$, and \ref{rule:E-App} then requires
$\fctx \types \fLt{\pntA}{\frm}$. $\tLink$ and $\tSpawn$ never block, so
their upper bound is $\pBot$. $\tNew$ puts both ends of a channel at one
point $\pviA$ and returns an existential, so a channel starts at a point
placed by its creation constraint and not by a global number.
$\tLink$ asks $\fEq{\pntA}{\pntB}$ because forwarding merges two
sessions into one conduit, and the actions of the merged channel need a
single position. Under $\fLt{\pntA}{\pntB}$ they would have two. The
condition is not an artefact of the static rule, \ref{rule:O-Link} of
the logical relation requires the world to alias the two points as well.
Whether a weaker condition that coerces one position into the other
would work is open.

\begin{figure}[tp]
  \adjustfigure[]
  \rulesection*{Type-indexed families of type signatures for constants $\cst : \ty$ (\rocq{konst_ty})}
  \begin{align*}
    & & & \text{side conditions the instance discharges}
    \\
    \tNew &{}: \tyUnit \lolli{\pBot}{\pTop}
      \tyEx{\pBot}{\pviA}{\cstr}{\lPair{\sAt{\sty}{\pviA}}{\sAt{\dual{\sty}}{\pviA}}}
    & & \cstr \text{ a creation constraint}
    \\
    \tSpawn &{}: (\tyUnit \lolli{\ubd}{\lbd} \tyUnit) \lolli{\pBot}{\pTop} \tyUnit
    \\
    \tSend &{}: \lPair{\tysub{\ty}{\pntA}{\pself}}{\sAt{(\sSend{\ty}{\sty})}{\pntA}}
      \lolli{\pntA}{\pTop} \sAt{\sty}{(\poff{\pntA}{1})}
    & & \fctx \types \fLt{\pntA}{\pnode} \text{ for every }
        \pnode \in \heads(\tysub{\ty}{\pntA}{\pself})
    \\
    \tRecv &{}: \sAt{(\sRecv{\ty}{\sty})}{\pntA}
      \lolli{\pntA}{\pTop} \lPair{\tysub{\ty}{\pntA}{\pself}}{\sAt{\sty}{(\poff{\pntA}{1})}}
    \\
    \tClos &{}: \sAt{\sTerm}{\pntA} \lolli{\pntA}{\pTop} \tyUnit
    \\
    \tLink &{}: \lPair{\sAt{\sty}{\pntA}}{\sAt{\dual{\sty}}{\pntB}}
      \lolli{\pBot}{\pTop} \tyUnit
    & & \fctx \types \fEq{\pntA}{\pntB}
    \\
    \tForce &{}: \tyBox{\pvivec}{\cstr}{\tyUnit \lolli{\ubd}{\pTop} \ty}
      \lolli{\tysub{\ubd}{\pntvec}{\pvivec}}{\pTop} \tysub{\ty}{\pntvec}{\pvivec}
    & & \fctx \types \tysub{\cstr}{\pntvec}{\pvivec}
  \end{align*}
  \rulesectionend
  \caption{Signatures of the constants}
  \Description{The seven type signatures, with the side conditions on
    points that each one reads from the fact theory.}
  \label{fig:signatures}
\end{figure}

\begin{definition}[Creation constraint \rocq{cc_ok}]
\label{def:creation}
  A point of $\pviA$ is a point with root $\pviA$. A constraint $\cstr$
  over one bound variable $\pviA$ is a \emph{creation constraint} under
  $\fctx$ when it is
  \begin{itemize}
    \item \emph{acyclic}:
      $\fctx, \cstr \nvdash \fLt{\pntA}{\pntA}$ for every point $\pntA$;
      and
    \item \emph{conservative}:
      $\fctx, \cstr \types \fLt{\pntA}{\pntB}$ implies
      $\fctx \types \fLt{\pntA}{\pntB}$, and
      $\fctx, \cstr \types \fEq{\pntA}{\pntB}$ implies
      $\fctx \types \fEq{\pntA}{\pntB}$, for all points
      $\pntA, \pntB$ without root $\pviA$.
  \end{itemize}
\end{definition}

A new channel may sit anywhere the context leaves room, and it may not
order or identify two points the context already has.
$\cstr = \set{\fEq{\pviA}{\pntA}}$ places the channel at $\pntA$, and
$\cstr = \set{\fLt{\ubd}{\pviA}, \fLt{\pviA}{\asTy{v}}}$ places it
between $\ubd$ and $\asTy{v}$, which conservativity allows only when
$\fctx \types \fLt{\ubd}{\asTy{v}}$. We write
$\tNewC{\pviA}{\cstr}$ for $\tNew$ at a fresh point $\pviA$ under
$\cstr$, $\tBoxA{\vr}{\pvivec}{\cstr}{\tm}$ for the box of
\ref{rule:E-Exp}, $\tAbsA{\ubd}{\lbd}{\vr}{\tm}$ for the abstraction of
\ref{rule:E-Abs}, and $\tForceA{\vr}{\pntvec}$ for the instantiation of
\ref{rule:E-Const} at $\tForce$. The derived form
$\tLetS{\tPair{\ch}{\chB}}{\tApp{\tNew}{\tUnit}}{\tm}$ unfolds to
\ref{rule:E-LetPair}, \ref{rule:E-Pair}, \ref{rule:E-App},
\ref{rule:E-Let*}, \ref{rule:E-Open} and \ref{rule:E-LetPair}
(\rocq{t_new_open}).

\paragraph{Processes}
\Cref{fig:procTypeSys} types processes. All return types are
\emph{primitive} (\rocq{prim}), so a thread cannot avoid closing its
channels by returning them. \ref{rule:P-Connect} composes two processes
that interact in a dual way on the channels of $\senv_{\asTy{C}}$ and
subsumes the mix rule at $\senv_{\asTy{C}} = \tEmpty$.

The process judgment $\senv \typesW{\wrld} \proc : \ty$ carries a world
$\wrld$, and so does the term judgment
$\tenv ; \lenv \typesFW{\frm}{\wrld} \tm : \ty$ that
\ref{rule:P-Term} reads (\rocq{typed}). Neither is the judgment
$\fctx ; \lenv \typesF{\frm} \tm : \ty$ of
\cref{fig:typeSys,fig:typeSysPoints}, which carries facts and no world.
\Cref{sec:worlds} defines worlds, \cref{sec:logicalRelation:relation}
defines the term environment $\tenv$, and \rocq{styped_typed} embeds
every source derivation into the world-carrying judgment at the empty
world. The restricted channels of \ref{rule:P-Term} are the domain of
$\tenv_{\asTy{B}}$, which holds both ends of each of them, so a thread
that owns a whole channel is typed without the channel appearing in its
interface.

\begin{figure}[tp]
  \adjustfigure[]
  \rulesection*{Primitive types $\primitive \ty$ (\rocq{prim})}
  \begin{proofrules}
    \infer*{}{
      \primitive \tyUnit
    }
    \and
    \infer*{
      \primitive \tyB \\ \primitive \tyC
    }{
      \primitive \lPair{\tyB}{\tyC}
    }
    \and
    \infer*{
      \primitive \tyB \\ \primitive \tyC
    }{
      \primitive \lSum{\tyB}{\tyC}
    }
    \and
    \infer*{
      \primitive \tyB
    }{
      \primitive \tyBox{\pvivec}{\cstr}{\tyUnit \lolli{\ubd}{\pTop} \tyB}
    }
  \end{proofrules}
  \rulesection{Typing rules for processes $\senv \typesW{\wrld} \proc : \ty$ (\rocq{ptyped})}
  \begin{proofrules}
    \infer*[lab=P-Term]{
      \primitive \ty \\
      \tenv_{\asTy{B}} \text{ holds both ends of each of its channels,
        dual at aliased points} \\
      \senv , \tenv_{\asTy{B}} ; \tEmpty \typesFW{\emptyset}{\wrld} \tm : \ty
    }{
      \senv \typesW{\wrld} (\dom(\tenv_{\asTy{B}}) , \mset{\tm}) : \ty
    } \label{rule:P-Term}
    \and
    \infer*[lab=P-Connect]{
      \senv_1 , \senv_{\asTy{C}} \typesW{\wrld} (\chset_1 , \pool_1) : \ty \\
      \senv_2 , \dual{\senv_{\asTy{C}}} \typesW{\wrld} (\chset_2 , \pool_2) : \tyUnit \\
      \dom(\senv_{\asTy{C}}) = \chset_0
    }{
      \senv_1 , \senv_2 \typesW{\wrld} (\chset_1 \cupdot \chset_2 \cupdot \chset_0 , \pool_1 \uplus \pool_2) : \ty
    } \label{rule:P-Connect}
  \end{proofrules}
  \rulesectionend
  \caption{Typing for processes}
  \Description{The four rules defining primitive types and the two
    process typing rules.}
  \label{fig:procTypeSys}
\end{figure}

\paragraph{Soundness}
A program starts under one constraint $\cstr_0$ over the top-level point
variables it begins with. A program that creates all its channels with
$\tNew$ takes $\cstr_0 = \emptyset$. A program that wants a global
order, such as the translation of a priority system, states it in
$\cstr_0$ and aliases its channels to the variables $\cstr_0$ speaks of.
The source programs of this paper name no point constant. Their points
are the variables that $\cstr_0$, a $\Box$ and an $\exists$ bind, and
\cref{thm:soundness-source} supplies the constants itself
(\rocq{c0_consts}). For such a program $\tysub{\ty}{\pntA}{\pself}$ is
the only way a payload acquires a point. A translated program may name
constants, and the image of PGV does (\rocq{proot}).

\begin{theorem}[Soundness for source programs \rocq{soundness_source}]
\label{thm:soundness-source}
  If $\cstr_0$ is acyclic, $\primitive \ty$ and
  $\cstr_0 ; \tEmpty \typesF{\emptyset} \tm : \ty$, then the thread
  $(\emptyset, \mset{\tm})$ is deadlock free.
\end{theorem}

\begin{proof}
  The proof builds the world the run needs out of $\cstr_0$. Embed the
  source derivation into the world-carrying judgment at the empty world,
  which supports no fact beyond equality and program order
  (\rocq{wnil}, \rocq{wnil_strict}, \rocq{styped_typed}). Let $n$ be the
  number of top-level point variables and send each of them to its own
  point constant. Acyclicity of $\cstr_0$ makes the world generated by
  the image of $\cstr_0$ under that map strict (\rocq{wc0},
  \rocq{wc0_strict}), it extends the empty world (\rocq{wc0_ext}), and
  the hypothesis $(\cstr_0, \pvivec)$ holds in it at those constants
  (\rocq{wc0_holds}). Instantiating the derivation along the map
  therefore yields the infinitary judgment in that world
  (\rocq{typed_sem_rules}, \rocq{c0_consts}), hence a semantic judgment
  (\rocq{sem_rules_fundamental}). One thread with nothing restricted is
  a process of the relation (\rocq{Prel_term}), and
  \cref{thm:adequacy} concludes. Point instantiation preserves
  $\primitive$ (\rocq{prim_pinst}), which \rocq{Prel_term} needs.
\end{proof}

Acyclicity of $\cstr_0$ is the acyclicity half of
\cref{def:creation}, checked against the empty context. The judgment of
\cref{thm:soundness-source} carries no world and no channel environment.
The world-carrying judgment $\senv \typesW{\wrld} \proc : \ty$ of
\cref{fig:procTypeSys} and the logical relation of
\cref{sec:logicalRelation} sit behind it. \rocq{styped_typed} bridges
them, taking a source derivation to the world-carrying term judgment at
any world and at the empty term environment, and \rocq{typed_styped}
comes back from a world with empty support.

\begin{example}[Delegation \rocq{deleg_main}]
\label{ex:delegation}
  One channel carries another, and the delegated session runs strictly
  inside the gap between the carrier's two actions. The two protocols
  are
  \begin{align*}
    \sty^0 &= \sSend{\tyUnit}{\sTerm}
    & \sty &= \sSend{(\sAt{\sty^0}{\pslot{\pself}{0}})}{\sTerm}
  \end{align*}
  and the program, with every annotation inlined, is
  \begin{align*}
    \asTmId{deleg} := {}
      & \tLetS{\tPair{\ch : \sAt{\sty}{\pviA}}{\ch' : \sAt{\dual{\sty}}{\pviA}}}{%
          \tApp{\tNewC{\pviA}{\fLt{\pviA}{\pslot{\pviA}{0}},\ %
            \fLt{\poff{\pslot{\pviA}{0}}{1}}{\poff{\pviA}{1}}}}{\tUnit}}{}
      \\
      & \tLetS{\tPair{\chB : \sAt{\sty^0}{\pslot{\pviA}{0}}}{%
          \chB' : \sAt{(\sRecv{\tyUnit}{\sTerm})}{\pslot{\pviA}{0}}}}{%
          \tApp{\tNewC{\pviB}{\fEq{\pviB}{\pslot{\pviA}{0}}}}{\tUnit}}{}
      \\
      & \tSeq{\tApp{\tSpawn}{\asTmId{K}}}{}
        \tLetS{\ch_1 : \sAt{\sTerm}{(\poff{\pviA}{1})}}{\tApp{\tSend}{\tPair{\chB}{\ch}}}{}
      \\
      & \tLet[\ivx][\chB_1]{\tApp{\tRecv}{\chB'}}{}
        \tSeq{\ivx}{\tSeq{\tApp{\tClos}{\chB_1}}{\tApp{\tClos}{\ch_1}}}
      \\
    \asTmId{K} := {}
      & \tAbsA{\poff{\pviA}{1}}{\pviA}{\tUnit}{}
        \tLet[\asTm{e}][\ch'_1]{\tApp{\tRecv}{\ch'}}{}
        \tSeq{\tApp{\tClos}{(\tApp{\tSend}{\tPair{\tUnit}{\asTm{e}}})}}{\tApp{\tClos}{\ch'_1}}
  \end{align*}
  The creation constraint of $\ch$ puts the delegated session strictly
  inside the gap between the carrier's two actions, and the second
  $\tNew$ creates $\chB$ at the carried root of a message the carrier
  has not sent yet. The order is one chain,
  \[
    \fLt{\pviA}{{}} \pslot{\pviA}{0}
    \prec \poff{\pslot{\pviA}{0}}{1}
    \prec \poff{\pviA}{1}
    \prec \pfv ,
  \]
  built from two facts of the first $\tNew$, program order, and the
  frame hypothesis \ref{rule:E-Abs} adds for the spawned thread. The
  rank $\poff{\pviA}{k} \mapsto 10k$ and
  $\poff{\pslot{\pviA}{0}}{o} \mapsto 5+o$ witnesses acyclicity.
  \Cref{sec:app-delegation} lists every blocking action with the premise
  it owes. The judgment is
  $\emptyset ; \tEmpty \typesF{\emptyset} \asTmId{deleg} : \tyUnit$
  (\rocq{deleg_typed}), so $\asTmId{deleg}$ is deadlock free
  (\rocq{deleg_deadlock_free}).

  A priority system cannot type this program. The carrier typed
  $\sSend{}{}^0.\sRecv{}{}^1$ leaves no integer between its two actions,
  and the number of actions the delegated session performs is not known
  where the carrier is typed.
\end{example}

%%% Local Variables:
%%% mode: latex
%%% TeX-master: "main"
%%% End:


\section{Semantic Session Typing for Deadlock Freedom}
\label{sec:logicalRelation}

The semantic typing judgment $\senv \proves \proc : \ty$ of this section
fulfils three goals.
\begin{enumerate}[(1)]
  \item It ensures deadlock freedom for closed processes (``adequacy'');
  \item
    It validates all rules of the static system of
    \cref{sec:syntacticTyping}, which gives the ``fundamental property''
    that
    $
      \senv \types \proc : \ty
      \implies
      \senv \proves \proc : \ty
    $;
  \item It contains safe but untypable processes.
\end{enumerate}


\subsection{Worlds and Readiness}
\label{sec:worlds}

The relation is indexed by a \emph{world}, which carries the order the
points of a running program stand in.

\begin{definition}[World \rocq{world}, \rocq{world_wf}, \rocq{strict}]
\label{def:world}
  A world $\wrld$ is a finite set of names, its \emph{support}, together
  with two relations $\fEq{}{}$ and $\fLt{}{}$ on points. It is
  \emph{well formed} when $\fEq{}{}$ is an equivalence closed under
  successor and slot congruence, $\fLt{}{}$ is transitive, closed under
  $\fEq{}{}$ on both sides and contains program order
  $\fLt{\pntA}{\poff{\pntA}{1}}$, every support name is a point
  constant, and outside the support only reflexivity and program order
  hold. It is \emph{strict} when $\fLt{}{}$ is irreflexive.
  $\wrld$ \emph{knows} $\pntA$ when $\rootof(\pntA)$ is in its support
  (\rocq{pt_known}).
\end{definition}

\begin{definition}[Extension \rocq{wext}, \rocq{wext_fresh}]
\label{def:wext}
  $\wrld \wextsym \wrld'$ holds when $\wrld'$ is strict and contains the
  support and both relations of $\wrld$.
  $\wrld \wextfresh \wrld'$ \emph{placing $N$} holds in addition when
  $N$ is disjoint from the support of $\wrld$, the support of $\wrld'$
  is that of $\wrld$ with $N$, and $\wrld'$ adds no fact between points
  whose roots avoid $N$.
\end{definition}

Only $\tNew$ extends the world, and it extends it freshly. The two
results below are what makes a composed run strict.

\begin{theorem}[Amalgamation \rocq{amalgamation}]
\label{thm:amalgamation}
  For strict $\wrld_1$ and $\wrld_2$, the join
  $\wrld_1 \sqcup \wrld_2$ is strict if and only if the join of their
  restrictions to the names both support is strict.
\end{theorem}

\begin{theorem}[Fresh extension \rocq{fresh_extension}]
\label{thm:fresh-extension}
  If $\wrld_1 \wextfresh \wrld_1'$ placing $N$, the support of
  $\wrld_2$ avoids $N$, $\wrld_2$ is strict and
  $\wrld_1 \sqcup \wrld_2$ is strict, then
  $\wrld_1' \sqcup \wrld_2$ is strict.
\end{theorem}

The step clauses of the relation are gated by \emph{readiness}. A term
or process with no internal step left must either be done or offer some
of the entries it holds, and the offered entries must lie strictly below
everything it holds but does not offer. Entries are compared, not names:
a thread holding both ends of $\ch$ and blocking on one of them holds
the other entry unoffered and would need
$\fLt{\heads(\ch)}{\heads(\ch)}$.

\begin{definition}[Readiness \rocq{ready_t}, \rocq{ready_p}]
\label{def:ready}
  $\readyop_{\wrld}(\tenv \mid \frm, \tm)$ holds when, provided $\tm$
  has no $\beta$, $\lNew{\ch}$ or $\lSpawnL{\val}$ step,
  \begin{enumerate*}[(i)]
    \item $\tm$ is a value if it offers nothing; and
    \item $\tm$ has a forwarding redex, or every entry it offers lies in
      $\wrld$ strictly below every entry of $\tenv$ it does not offer
      and strictly below every point of $\frm$.
  \end{enumerate*}
  $\readyop_{\wrld}(\senv, \proc)$ holds when, provided $\proc$ has no
  $\tau$ step, $\proc$ has a forwarding redex, or $\proc$ is terminated
  if it offers nothing and every entry it holds and does not offer lies
  strictly above some entry it offers.
\end{definition}

A forwarding never blocks, so a term or process at a forwarding redex
compares nothing.

\subsection{A Logical Relation for \OurLang}
\label{sec:logicalRelation:relation}

Three relations express the typing constraints
(\cref{fig:logrel-rules-value,fig:logrel-rules-term,fig:logrel-rules-proc}):
$\lrVW{\wrld}{n}{\tenv \types \ty}$ for values,
$\lrEW{\wrld}{n}{\tenv \types \ty}{\frm}$ for terms, and
$\lrPW{\wrld}{n}{\senv \types \ty}$ for processes.
The index $n \in \Nat$ is a step index~\cite{Ahmed06} bounding the
number of transitions the relation examines.
A session environment $\senv$ maps each channel to one
$\sAt{\sty}{\pntA}$, and a term environment $\tenv$ maps a channel to
one or both of its ends, because a term may hold both
(\rocq{senv}, \rocq{tenv}). $\heads(\tenv)$ is the set of points of the
entries of $\tenv$, and $\tenv_1 \cdot \tenv_2$ splits $\tenv$ up to
$\fEq{}{}$ of the world (\rocq{tenv_comp}).

\begin{figure}[tp]
  \adjustfigure[]
  \rulesection*{Value logical relation $\val \in \lrVW{\wrld}{n}{\tenv \types \ty}$ (\rocq{Vdef})}
  \begin{proofrules}
    \infer*[lab=V-Unit]{}{
      \tUnit \in \lrVW{\wrld}{n}{\tEmpty \types \tyUnit}
    } \label{rule:V-Unit}
    \and
    \infer*[lab=V-Chan]{
      \sty_0 \approx \sty \\
      \wrld \types \fEq{\pntB}{\pntA}
    }{
      \ch \in \lrVW{\wrld}{n}{\ch : \sAt{\sty_0}{\pntB} \types \sAt{\sty}{\pntA}}
    } \label{rule:V-Chan}
    \and
    \infer*[lab=V-Const]{
      \cst : \ty_0 \text{ an instance of a signature under } \wrld \\
      \ty_0 \approx_{\wrld} \ty
    }{
      \cst \in \lrVW{\wrld}{n}{\tEmpty \types \ty}
    } \label{rule:V-Const}
    \and
    \infer*[lab=V-Pair]{
      \val_1 \in \lrVW{\wrld}{n}{\tenv_1 \types \ty} \\
      \val_2 \in \lrVW{\wrld}{n}{\tenv_2 \types \tyB}
    }{
      \tPair{\val_1}{\val_2} \in \lrVW{\wrld}{n}{\tenv_1 \cdot \tenv_2 \types \lPair{\ty}{\tyB}}
    } \label{rule:V-Pair}
    \and
    \infer*[lab=V-Inj]{
      \val \in \lrVW{\wrld}{n}{\tenv \types \ty_{\asTy{i}}}
    }{
      \tInj{\asTm{i}}{\val} \in \lrVW{\wrld}{n}{\tenv \types \lSum{\ty_1}{\ty_2}}
    } \label{rule:V-Inj}
    \and
    \infer*[lab=V-Abs]{
      \fchans(\tm) = \dom(\tenv) \\
      \wrld \types \fLe{\lbd}{\heads(\tenv)}
      \\\\
      \forall \wrld' \wextsym \wrld,\ j < n,\ \tenv'',\
        \val \in \lrVW{\wrld'}{j}{\tenv' \types \ty},\
        \frm \text{ with } \wrld' \types \fLt{\ubd}{\frm}
      \\
      \text{and } \tenv'' = \tenv \cdot \tenv' \st
        \tSub{\tm}{\val}{\vr} \in \lrEW{\wrld'}{j}{\tenv'' \types \tyB}{\frm}
    }{
      \tAbs{\vr}{\tm} \in \lrVW{\wrld}{n}{\tenv \types \ty \lolli{\ubd}{\lbd} \tyB}
    } \label{rule:V-Abs}
    \and
    \infer*[lab=V-Box]{
      \fchans(\tm) = \emptyset
      \\\\
      \forall \wrld' \wextsym \wrld,\ \pntvec \text{ known to } \wrld'
        \text{ with } \wrld' \types \tysub{\cstr}{\pntvec}{\pvivec} \st
      \\
      \tAbs{\tUnit}{\tSub{\tm}{\tBoxP{\vr}{\tm}}{\vr}} \in
        \lrVW{\wrld'}{n}{\tEmpty \types \tysub{(\tyUnit \lolli{\ubd}{\pTop} \ty)}{\pntvec}{\pvivec}}
    }{
      \tBoxP{\vr}{\tm} \in \lrVW{\wrld}{n}{\tEmpty \types \tyBox{\pvivec}{\cstr}{\tyUnit \lolli{\ubd}{\pTop} \ty}}
    } \label{rule:V-Box}
    \and
    \infer*[lab=V-Ex]{
      \wrld \types \fLe{\lbd}{\heads(\tenv)} \\
      \pntvec \text{ known to } \wrld \\
      \wrld \types \tysub{\cstr}{\pntvec}{\pvivec} \\
      \val \in \lrVW{\wrld}{n}{\tenv \types \tysub{\ty}{\pntvec}{\pvivec}}
    }{
      \val \in \lrVW{\wrld}{n}{\tenv \types \tyEx{\lbd}{\pvivec}{\cstr}{\ty}}
    } \label{rule:V-Ex}
  \end{proofrules}
  \rulesectionend
  \caption{The value logical relation}
  \Description{The seven rules of the value relation.}
  \label{fig:logrel-rules-value}
\end{figure}

\begin{figure}[tp]
  \adjustfigure[]
  \rulesection*{Term logical relation $\tm \in \lrEW{\wrld}{n}{\tenv \types \ty}{\frm}$ (\rocq{Edef})}
  \begin{proofrules}
    \infer*[lab=E-Zero]{}{
      \tm \in \lrEW{\wrld}{0}{\tenv \types \ty}{\frm}
    } \label{rule:E-Zero}
    \and
    \infer*[lab=E-Intro]{
      \irred \tm \implies \tm \in \lrVW{\wrld}{n+1}{\tenv \types \ty}
      \\
      \neg\irred \tm \implies \readyop_{\wrld}(\tenv \mid \frm, \tm)
      \\
      \forall \lambda, \tm'\st \tm \trans{\lambda} \tm' \implies
        \obl{\wrld ; \tenv \mid \frm}{n+1}{\tm \trans{\lambda} \tm' : \ty}
    }{
      \tm \in \lrEW{\wrld}{n+1}{\tenv \types \ty}{\frm}
    } \label{rule:E-Intro}
  \end{proofrules}
  \rulesection{Obligations of a term transition $\obl{\wrld ; \tenv \mid \frm}{n}{\tm \trans{\lambda} \tm' : \ty}$ (\rocq{Eobl})}
  \begin{proofrules}
    \infer*[lab=O-Beta]{
      \tm' \in \lrEW{\wrld}{n \cost{\lambda}}{\tenv \types \ty}{\frm}
    }{
      \obl{\wrld ; \tenv \mid \frm}{n}{\tm \trans{\beta} \tm' : \ty}
    } \label{rule:O-Beta}
    \and
    \infer*[lab=O-New]{
      \ch \notin \dom(\tenv) \implies
      \exists\, \pcst \text{ fresh for } \wrld, \tenv, \ty, \frm,\
        \sty,\ \wrld \wextfresh \wrld' \text{ placing } \pcst \st
      \\
      \tm' \in \lrEW{\wrld'}{n \cost{\lambda}}{\tenv , \ch : \sAt{\sty}{\pcst} \cdot \sAt{\dual{\sty}}{\pcst} \types \ty}{\frm}
    }{
      \obl{\wrld ; \tenv \mid \frm}{n}{\tm \trans{\lNew{\ch}} \tm' : \ty}
    } \label{rule:O-New}
    \and
    \infer*[lab=O-Spawn]{
      \val \in \lrVW{\wrld}{n}{\tenv_1 \types \tyUnit \lolli{\ubd}{\lbd} \tyUnit} \\
      \tm' \in \lrEW{\wrld}{n \cost{\lambda}}{\tenv_2 \types \ty}{\frm}
    }{
      \obl{\wrld ; \tenv_1 \cdot \tenv_2 \mid \frm}{n}{\tm \trans{\lSpawnL{\val}} \tm' : \ty}
    } \label{rule:O-Spawn}
    \and
    \infer*[lab=O-Recv]{
      \forall \wrld' \wextsym \wrld,\ j \le n \cost{\lambda},\ \tenv_2,\
        \val \in \lrVW{\wrld'}{j}{\tenv_2 \types \tysub{\tyB}{\pntA}{\pself}} \st
      \\
      \tm' \in \lrEW{\wrld'}{j}{(\tenv_1 , \ch : \sAt{\sty}{(\poff{\pntA}{1})}) \cdot \tenv_2 \types \ty}{\frm}
    }{
      \obl{\wrld ; \tenv_1 , \ch : \sAt{(\sRecv{\tyB}{\sty})}{\pntA} \mid \frm}{n}{\tm \trans{\lRecv{\ch}{\val}} \tm' : \ty}
    } \label{rule:O-Recv}
    \and
    \infer*[lab=O-Send]{
      \val \in \lrVW{\wrld}{n}{\tenv_2 \types \tysub{\tyB}{\pntA}{\pself}} \\
      \tm' \in \lrEW{\wrld}{n \cost{\lambda}}{\tenv_1 , \ch : \sAt{\sty}{(\poff{\pntA}{1})} \types \ty}{\frm}
    }{
      \obl{\wrld ; (\tenv_1 , \ch : \sAt{(\sSend{\tyB}{\sty})}{\pntA}) \cdot \tenv_2 \mid \frm}{n}{\tm \trans{\lSend{\ch}{\val}} \tm' : \ty}
    } \label{rule:O-Send}
    \and
    \infer*[lab=O-Close]{
      \tm' \in \lrEW{\wrld}{n \cost{\lambda}}{\tenv_1 \types \ty}{\frm}
    }{
      \obl{\wrld ; \tenv_1 , \ch : \sAt{\sTerm}{\pntA} \mid \frm}{n}{\tm \trans{\lTerm{\ch}} \tm' : \ty}
    } \label{rule:O-Close}
    \and
    \infer*[lab=O-Link]{
      \sty_{\chB} \isdual \sty_{\ch} \\
      \wrld \types \fEq{\pntA}{\pntB} \\
      \tm' \in \lrEW{\wrld}{n \cost{\lambda}}{\tenv_1 \types \ty}{\frm}
    }{
      \obl{\wrld ; \tenv_1 , \ch : \sAt{\sty_{\ch}}{\pntA} , \chB : \sAt{\sty_{\chB}}{\pntB} \mid \frm}{n}{}
      \\ \qquad {\tm \trans{\lLink{\ch}{\chB}} \tm' : \ty}
    } \label{rule:O-Link}
  \end{proofrules}
  \rulesectionend
  \caption{The term logical relation}
  \Description{The two rules of the term relation and the seven
    obligation rules of a term transition.}
  \label{fig:logrel-rules-term}
\end{figure}

A send, a close and a forwarding give a resource away, so they keep the
index. $n \cost{\lambda} = n$ for those three labels and
$n \cost{\lambda} = n - 1$ for every other (\rocq{tcost}). The
corresponding successor term is strictly smaller, which is what makes
\cref{fig:logrel-rules-term} a definition (\rocq{Efuel}).
\Cref{rule:V-Abs} reads the argument and the result at strictly smaller
indices only, so \ref{rule:V-Box} may use it at an instance of its own
type.

\begin{remark}
\label{rm:vabs}
  The side conditions on free channels of \ref{rule:V-Abs} and
  \ref{rule:V-Box} are needed because the LTS decides which channels a
  send extrudes by looking at the channels the sent value mentions
  syntactically (\cref{rm:sentc}). They are the only syntactic
  conditions in the relation, and they rule out a sent value that
  mentions channels in dead code which the continuation of the send
  still uses.
\end{remark}

\begin{figure}[tp]
  \adjustfigure[]
  \rulesection*{Process logical relation $\proc \in \lrPW{\wrld}{n}{\senv \types \ty}$ (\rocq{Pdef})}
  \begin{proofrules}
    \infer*[lab=P-Zero]{}{
      \proc \in \lrPW{\wrld}{0}{\senv \types \ty}
    } \label{rule:P-Zero}
    \and
    \infer*[lab=P-Intro]{
      \irred \proc \implies \proc = (\asPrOp{\emptyset}, \mset{\val,\tUnit,\dots,\tUnit})
        \text{ with } \val \in \lrVW{\wrld}{n+1}{\senv \types \ty}
      \\
      \neg\irred \proc \implies \readyop_{\wrld}(\senv, \proc)
      \\
      \forall \pi, \proc'\st \proc \trans{\pi} \proc' \implies
        \obl{\wrld ; \senv}{n+1}{\proc \trans{\pi} \proc' : \ty}
    }{
      \proc \in \lrPW{\wrld}{n+1}{\senv \types \ty}
    } \label{rule:P-Intro}
  \end{proofrules}
  \rulesection{Obligations of a process transition $\obl{\wrld ; \senv}{n}{\proc \trans{\pi} \proc' : \ty}$ (\rocq{Pobl})}
  \begin{proofrules}
    \infer*[lab=Q-Tau]{
      \proc'.1 \disjoint \dom(\senv) \implies
      \exists\, N \text{ fresh for } \senv, \ty,\
        \wrld \wextfresh \wrld' \text{ placing } N \st
      \proc' \in \lrPW{\wrld'}{n \cost{\pi}}{\senv \types \ty}
      \\
      \text{and, if the step is the $\tau$ of \ref{rule:P-Act-Link}},\
        \proc' \in \lrPW{\wrld}{n}{\senv \types \ty}
    }{
      \obl{\wrld ; \senv}{n}{\proc \trans{\tau} \proc' : \ty}
    } \label{rule:Q-Tau}
    \and
    \infer*[lab=Q-Recv]{
      \forall \wrld' \wextsym \wrld,\ j \le n \cost{\pi},\ \tenv_2
        \text{ with } \dom(\tenv_2) \disjoint \proc.1,\
        \val \in \lrVW{\wrld'}{j}{\tenv_2 \types \tysub{\tyB}{\pntA}{\pself}} \st
      \\
      (\proc'.1 \cup B, \proc'.2) \in
        \lrPW{\wrld'}{j}{(\senv_1 , \ch : \sAt{\sty}{(\poff{\pntA}{1})}) \mathbin{\triangleright} \tenv_2 \types \ty}
    }{
      \obl{\wrld ; \senv_1 , \ch : \sAt{(\sRecv{\tyB}{\sty})}{\pntA}}{n}{\proc \trans{\lRecv{\ch}{\val}} \proc' : \ty}
    } \label{rule:Q-Recv}
    \and
    \infer*[lab=Q-Send]{
      \dom(\tenv_{\asTy{e}}) = \fchans(\val) \cap \proc.1 \\
      \tenv_{\asTy{e}} \text{ and } \senv_{\asTy{e}} \text{ are dual at aliased points}
      \\\\
      \val \in \lrVW{\wrld}{n}{\senv_{\asTy{v}} , \tenv_{\asTy{e}} \types \tysub{\tyB}{\pntA}{\pself}}
      \quad
      \proc' \in \lrPW{\wrld}{n \cost{\pi}}{\senv_0 , \senv_{\asTy{e}} , \ch : \sAt{\sty}{(\poff{\pntA}{1})} \types \ty}
    }{
      \obl{\wrld ; \senv_0 , \senv_{\asTy{v}} , \ch : \sAt{(\sSend{\tyB}{\sty})}{\pntA}}{n}{\proc \trans{\lSend{\ch}{\val}} \proc' : \ty}
    } \label{rule:Q-Send}
    \and
    \infer*[lab=Q-Close]{
      \proc' \in \lrPW{\wrld}{n \cost{\pi}}{\senv_1 \types \ty}
    }{
      \obl{\wrld ; \senv_1 , \ch : \sAt{\sTerm}{\pntA}}{n}{\proc \trans{\lTerm{\ch}} \proc' : \ty}
    } \label{rule:Q-Close}
    \and
    \infer*[lab=Q-Link]{
      \ch \ne \chB \\
      \sty_{\chB} \isdual \sty_{\ch} \\
      \wrld \types \fEq{\pntA}{\pntB} \\
      \proc' \in \lrPW{\wrld}{n \cost{\pi}}{\senv_1 \types \ty}
    }{
      \obl{\wrld ; \senv_1 , \ch : \sAt{\sty_{\ch}}{\pntA} , \chB : \sAt{\sty_{\chB}}{\pntB}}{n}{\proc \trans{\lLink{\ch}{\chB}} \proc' : \ty}
    } \label{rule:Q-Link}
  \end{proofrules}
  \rulesectionend
  \caption{The process logical relation}
  \Description{The two rules of the process relation and the five
    obligation rules of a process transition.}
  \label{fig:logrel-rules-proc}
\end{figure}

\Cref{rule:Q-Tau} is the only clause that grows the world, by the fresh
extension that \ref{rule:P-Act-New} needs; a $\tau$ of
\ref{rule:P-Act-Link} creates no point and keeps the world, and it also
gives away a resource, so it carries its obligation at $n$ as well.
In \ref{rule:Q-Recv} the message may bring the other end of a channel
the receiver already holds, or both ends of a fresh one. Those channels
$B$ become restricted in the continuation and their entries leave both
environments, which the operation $\senv \mathbin{\triangleright} \tenv_2$
performs (\rocq{recv_back}, \rocq{recv_env}).
\Cref{rule:Q-Send} is where a private channel becomes part of the
external interface. The restricted channels the message carries are
$\dom(\tenv_{\asTy{e}}) = \fchans(\val) \cap \chset$; the ends that
leave with the message are typed in $\tenv_{\asTy{e}}$, the ends that
stay in $\senv_{\asTy{e}}$, and the two are dual at aliased points
(\rocq{send_split}). The continuation then serves
$\senv_{\asTy{e}}$ as part of its interface. No condition on the
topology of the continuation is imposed: the order on points does that
work, so this rule needs no split of the continuation into one process
per extruded channel.

The syntactic side condition $\fchans(\val)$ in \ref{rule:Q-Send} is
again a consequence of the design of the LTS (\cref{rm:sentc}).

\subsection{Fundamental Property}
\label{sec:fundamental}

The term relation applies to closed terms, so an open term is closed by
a grounding substitution.

\begin{definition}[Grounding substitutions \rocq{grounding}]
\label{def:grounding}
  \[
    \lrGW{\wrld}{n}{\tenv ; \vr_1 : \ty_1, \dots, \vr_k : \ty_k} :=
      \left\{
        [\vr_1/\val_1,\dots,\vr_k/\val_k] \;\middle|\;
          \begin{matrix}
          \tenv = \tenv_1 \cdot {} \dots \cdot \tenv_k
            \\
          \forall i.\,
            \val_i \in \lrVW{\wrld}{n}{\tenv_i \types \ty_i}
          \end{matrix}
      \right\}
  \]
\end{definition}

\begin{definition}[Semantic judgments \rocq{sem_typed}, \rocq{sem_ptyped}]
\label{def:sem-judgments}
  The judgment $\tenv ; \lenv \provesF{\frm}_{\wrld} \tm : \ty$ holds if
  \[
    \forall n \in \Nat.\,
      \forall \wrld' \wextsym \wrld.\,
        \forall \tenv'', \gamma \in \lrGW{\wrld'}{n}{\tenv' ; \lenv}
          \text{ with } \tenv'' = \tenv \cdot \tenv'.\,
          \tm\gamma \in \lrEW{\wrld'}{n}{\tenv'' \types \ty}{\frm}
  \]
  The judgment $\senv \proves_{\wrld} \proc : \ty$ holds if
  $
    \forall n \in \Nat.\,
      \proc \in \lrPW{\wrld}{n}{\senv \types \ty}
  $.
\end{definition}

The judgment is purely semantic. It says nothing about the syntax of
$\tm$ beyond what the term relation says about its behaviour. The two
syntactic conditions of \cref{rm:vabs} live in a separate predicate
$\mathit{wf}(\tenv; \lenv, \tm)$ (\rocq{term_wf}), which holds of every
syntactically typed term (\rocq{typed_term_wf_s}) and enters as a
premise of exactly the two compositionality lemmas that build a
suspended computation.

\begin{lemma}[\rocq{Prel_steps}]
\label{lm:lrproc-inv}
  If $\proc \in \lrPW{\wrld}{n+m}{\tEmpty \types \ty}$ and
  $\proc \redu[n] \proc'$, then
  $\proc' \in \lrPW{\wrld'}{m}{\tEmpty \types \ty}$ for some $\wrld'$.
\end{lemma}

\begin{lemma}[\rocq{Prel_safe}]
\label{lm:lrproc-safe}
  If $\proc \in \lrPW{\wrld}{n+1}{\tEmpty \types \ty}$ then $\proc$ is
  safe.
\end{lemma}

\begin{theorem}[Adequacy \rocq{adequacy}]
  \label{thm:adequacy}
  If $\tEmpty \proves_{\wrld} \proc : \ty$ then $\proc$ is deadlock
  free.
\end{theorem}

\begin{proof}
  Let $\proc \redu[m] \proc'$. Pick $n = m+1$ in
  $\tEmpty \proves_{\wrld} \proc : \ty$. By
  \cref{lm:lrproc-inv}, $\proc' \in \lrPW{\wrld'}{1}{\tEmpty \types \ty}$
  for some $\wrld'$, and by \cref{lm:lrproc-safe} $\proc'$ is safe.
\end{proof}

The relation validates every rule of the static system. Replacing
$\types$ by $\proves$ in \cref{fig:typeSys,fig:typeSysPoints,fig:procTypeSys} and reading
each inference rule as an implication gives a valid statement about the
semantic judgment, with the two syntactic premises noted above.
\Cref{tab:compositionality} names the declaration for each rule, and the
three lemmas below spell out the rules that carry a side condition of
their own.

\begin{table}[tp]
  \caption{Compositionality: one declaration per rule of the static
    system.}
  \label{tab:compositionality}
  \begin{tabular}{llllll}
    \toprule
    Rule & Declaration & Rule & Declaration & Rule & Declaration \\
    \midrule
    \ref{rule:E-Unit} & \rocq{sem_typed_unit}
      & \ref{rule:E-Const} & \rocq{sem_typed_konst}
      & \ref{rule:E-Var} & \rocq{sem_typed_var} \\
    \ref{rule:E-Abs} & \rocq{sem_typed_lam}
      & \ref{rule:E-App} & \rocq{sem_typed_app}
      & \ref{rule:E-Pair} & \rocq{sem_typed_pair} \\
    \ref{rule:E-LetPair} & \rocq{sem_typed_letp}
      & \ref{rule:E-Let*} & \rocq{sem_typed_letu}
      & \ref{rule:E-Inj} & \rocq{sem_typed_inj1} \\
    \ref{rule:E-Case} & \rocq{sem_typed_case}
      & \ref{rule:E-Weak} & \rocq{sem_typed_wk}
      & \ref{rule:E-Contract} & \rocq{sem_typed_contract} \\
    \ref{rule:E-Conv} & \rocq{sem_typed_teq}
      & \ref{rule:E-Exp} & \rocq{sem_typed_box}
      & \ref{rule:E-Pack} & \rocq{sem_typed_pack} \\
    \ref{rule:E-Open} & \rocq{sem_typed_open}
      & \ref{rule:P-Term} & \rocq{Prel_term}
      & \ref{rule:P-Connect} & \rocq{Prel_connect} \\
    \bottomrule
  \end{tabular}
\end{table}

A channel constant is typed by its own rule in the world-carrying
judgment, with declaration \rocq{sem_typed_chan}, and
$\tApp{\tForce}{\tm}$ is \ref{rule:E-App} at $\tForce$
(\rocq{sem_typed_force}). Composing two processes that do not interact
is \ref{rule:P-Connect} with no shared channel (\rocq{Prel_indep}).

\begin{lemma}[Compositionality: \ref{rule:E-Abs} \rocq{sem_typed_lam}]
  \label{lem:comp:abs}
  If $\tenv ; \lenv , \vr : \ty \provesF{\set{\pfv}}_{\wrld} \tm : \tyB$
  and $\wrld \types \fLe{\lbd}{\heads(\tenv) \cup \heads(\lenv)}$ and
  $\mathit{wf}(\tenv ; \lenv , \vr : \ty, \tm)$, then
  $\tenv ; \lenv \provesF{\frm}_{\wrld} \tAbs{\vr}{\tm}
    : \ty \lolli{\ubd}{\lbd} \tyB$.
\end{lemma}

\begin{lemma}[Compositionality: \ref{rule:E-Exp} \rocq{sem_typed_box}]
  \label{lem:comp:exp}
  If
  $\tEmpty ; \lenv , \vr : \tyBox{\pvivec}{\cstr}{\tyUnit \lolli{\ubd}{\pTop} \ty}
     \provesF{\set{\pfv}}_{\wrld} \tm : \ty$
  with $\lenv$ of $\Box$ types only and
  $\mathit{wf}(\tEmpty ; \lenv , \vr : \tyBox{\pvivec}{\cstr}{\tyUnit \lolli{\ubd}{\pTop} \ty}, \tm)$,
  then
  $\tEmpty ; \lenv \provesF{\frm}_{\wrld} \tBoxP{\vr}{\tm}
    : \tyBox{\pvivec}{\cstr}{\tyUnit \lolli{\ubd}{\pTop} \ty}$.
\end{lemma}

\begin{lemma}[Compositionality: \ref{rule:P-Connect} \rocq{Prel_connect}]
  \label{lem:comp:connect}
  If $\primitive \ty$, the two interfaces
  $\senv_{\asTy{C}}$ and $\dual{\senv_{\asTy{C}}}$ are dual at aliased
  points of $\wrld$,
  $(\chset_1, \pool_1) \in \lrPW{\wrld}{n}{\senv_1 , \senv_{\asTy{C}} \types \ty}$,
  $(\chset_2, \pool_2) \in \lrPW{\wrld}{n}{\senv_2 , \dual{\senv_{\asTy{C}}} \types \tyUnit}$,
  the restricted channels of each component avoid the interface of the
  other, and $\wrld \sqcup \wrld$ is strict, then
  \[
    (\chset_1 \cupdot \chset_2 \cupdot \dom(\senv_{\asTy{C}}),
     \pool_1 \uplus \pool_2)
      \in \lrPW{\wrld}{n}{\senv_1 , \senv_2 \types \ty}
  \]
\end{lemma}

\Cref{lem:comp:connect} is the linear-logic cut, and it is the one place
where strictness of the composed world is needed:
\cref{thm:amalgamation,thm:fresh-extension} supply it along the run,
because the only clause that grows a world grows it freshly
(\ref{rule:Q-Tau}) and the names it adds are unknown to the other
component.

\begin{theorem}[Fundamental property \rocq{fundamental}, \rocq{fundamental_proc}]
\label{th:fundamental}
  Typing implies semantic typing, at every strict world.
  \begin{itemize}
    \item
      $
        \fctx ; \lenv \typesF{\frm} \tm : \ty
        \implies
        \tenv ; \lenv \provesF{\frm}_{\wrld} \tm : \ty
      $
    \item
      $
        \senv \types \proc : \ty
        \implies
        \senv \proves_{\wrld} \proc : \ty
      $
  \end{itemize}
\end{theorem}

\begin{theorem}[Soundness \rocq{soundness}]
\label{th:soundness}
  If $\wrld$ is strict and $\tEmpty \types \proc : \ty$, then $\proc$ is
  deadlock free.
\end{theorem}

\Cref{thm:soundness-source} is \cref{th:soundness} for a source program:
the judgment of \cref{fig:typeSys} and \cref{fig:typeSysPoints} reads no world, the acyclic
$\cstr_0$ supplies a strict world that realizes it, and
\rocq{styped_typed} embeds the derivation (\rocq{soundness_source}).

%%% Local Variables:
%%% mode: latex
%%% TeX-master: "main"
%%% End:


\section{Extensions}
\label{sec:extensions}

In this section, we discuss some extensions which are all incorporated
in our proved mechanization. In particular, we cover sum types, replication and
recursion, recursive types, and the forwarding operator link. That is, all
features of GV \cite{conf/esop/LindleyM15} are covered.

\subsection{Sums}
\label{sec:extension:-sums}


\begin{figure}[tp]
  New values, terms, evaluation contexts \hfill \rocq{TSum}
  \begin{align*}
    \fFunc{v} &::= \dots \sepr \fFunc{\tInj{1}{v}} \sepr \fFunc{\tInj{2}{v}}
    & \text{injections}\\
    \fFunc{e} &::= \dots \sepr \fFunc{\tCase{e}{x}{e}{x}{e}} &
                                                                       \text{sum elimination}\\
    \fFunc{\fCtx{E}} & ::= \dots \sepr \fFunc{\tInj{1}{\fCtx{E}}} \sepr \fFunc{\tInj{2}{\fCtx{E}}}
  \end{align*}
  New labeled transitions \hfill \rocq{HeadCase1}
  \begin{mathpar}
    \begin{bussproof}[T-Act-Case]
      \bussAx{\fFunc{\tCase{(\tInj{i}{v})}{x_1}{e_1}{x_2}{e_2}} \trans{\beta} \fFunc{\tSub{e_i}{v}{x_i}}}
    \end{bussproof}
    % \and
    % \begin{bussproof}[T-Act-Case2]
    %   \bussAx{\fFunc{\tCase{\tInj{2}{v}}{\tLVA}{e}{\tLVA}{e}} \trans{\beta} \fFunc{\tSub{e'}{\tFixI{\tIVY}{e}}{\tIVX}}}
    % \end{bussproof}
  \end{mathpar}
  New typing rules for terms \hfill \rocq{typed_inj1}
  \begin{mathpar}
    \begin{bussproof}[E-Inj]
      \bussAs{\fType{\Sigma ; \Delta \vdash \fFunc{e} : T_i}}
      \bussU{\fType{\Sigma ; \Delta \vdash \fFunc{\tInj{i}{e}} : \lSum{T_1}{T_2}}}
    \end{bussproof}
    \and
    \begin{bussproof}[E-Case]
      \bussAs{\fType{\Sigma_1 ; \Delta_1 \vdash \fFunc{e} : \lSum{T_1}{T_2}}}
      \bussAsBelow{\fType{\Sigma_2 ; \Delta_2, \fFunc{x_1}:T_1 \vdash \fFunc{e_1} :U}}
      \bussAsBelow{\fType{\Sigma_2 ; \Delta_2, \fFunc{x_2}:T_2 \vdash \fFunc{e_2} :U}}
      \bussU{\fType{\Sigma_1 , \Sigma_2 ; \Delta_1 , \Delta_2 \vdash \fFunc{\tCase{e}{x_1}{e_1}{x_2}{e_2}} : U}}
    \end{bussproof}
  \end{mathpar}
  New case for the value logical relation \hfill \rocq{Vdef}
  \begin{align*}
    \lrVal{n}{\fType{\Sigma \vdash \lSum{T_1}{T_2}}}
  &= \set{ \fFunc{\tInj{1}{v}}  \mid
    \fFunc{v} \in \lrVal{n}{\fType{\Sigma \vdash T_1}}
    } \cup \set{ \fFunc{\tInj{2}{v}}  \mid
    \fFunc{v} \in \lrVal{n}{\fType{\Sigma \vdash T_2}}
    }
  \end{align*}
  New rule for primitive types \hfill \rocq{prim_sum}
  \begin{mathpar}
    \begin{bussproof}
      \bussAs{\primitive \fType{T}}
      \bussAs{\primitive \fType{U}}
      \bussB{\primitive \fType{\lSum{U}{V}}}
    \end{bussproof}
  \end{mathpar}
  \caption{Extensions for sum types}
  \Description{Additional values, terms and evaluation contexts for sums; one reduction axiom; two typing rules; one case of the value logical relation; and one rule for primitive types.}
  \label{fig:sum-types}
\end{figure}

Moving towards a realistic calculus, we wish to support branching
and choice in session types. It is well-known that branch and
choice can be simulated using higher-order sessions and sum types
\cite{DBLP:conf/ppdp/DardhaGS12,conf/esop/LindleyM15} by setting
\begin{align*}
  \fType{S_1 \ichoice S_2} &:= \fType{\sSend{(\dual{S_1} + \dual{S_2})}{\sTerm}}
  & \fType{S_1 \echoice S_2} &:= \fType{\sRecv{(S_1 + S_2)}{\sWait}}
\end{align*}
and encoding the operations to select and match on a channel as follows.
\begin{align*}
  \fFunc{\mathit{select}~\textit{left}~c} &:= \fFunc{\tFork (\lambda
                                   d. \tApp{\tClos}{(\tSend{\tPair{\tInj{1}{d}}{c}})})}
  &  \fFunc{\mathit{select}~\textit{right}~c} &:= \fFunc{\tFork (\lambda
                                       d. \tApp{\tClos}{(\tSend{\tPair{\tInj{2}{d}}{c}})})}
\end{align*}
\begin{align*}
  \fFunc{\mathit{match}~c~(x_1.e_1) (x_2.e_2)} &= \fFunc{\tLet[d][c]{\tRecv c}{\tClos{c}; \tCase{d}{x_1}{e_1}{x_2}{e_2}}}
\end{align*}
Directly supporting branch and choice would require additional labeled
reductions and hence additional cases in the definition of all logical
relations. To support the simulated branch and choice, it is sufficient to
implement sum types. To this end, we extend the syntax, labeled reduction,
typing, and the value logical relation in the usual way as
shown in \cref{fig:sum-types}. We prove the following compositionality lemmas
corresponding to the typing rules for sum types and integrate them in the proof
of the fundamental property. 

\begin{lemma}[Compositionality: \ruleLabel{inj} \rocq{sem_typed_inj1}]
  \label{lem:comp:inj}
  If
  $ \fType{\Sigma;\Delta \proves \fFunc{e} : T_i}$,
  then
  $ \fType{\Sigma;\Delta \proves \fFunc{\tInj{i}{e}} : \lSum{T_1}{T_2}}$,
  for $i\in\{1,2\}$.
\end{lemma}

\begin{lemma}[Compositionality: \ruleLabel{case} \rocq{sem_typed_case}]
  \label{lem:comp:case}
  If  $ \fType{\Sigma_1;\Delta_1 \proves \fFunc{e} : \lSum{T_1}{T_2}}$ and
  \begin{itemize}
  \item for all $ \fType{\Sigma;\Delta \proves \fFunc{v_1} : {T_1}}$ implies
    $ \fType{\Sigma,\Sigma_2;\Delta,\Delta_2 \proves \fFunc{e_1[v_1/x_1]} : U}$,
  \item for all $ \fType{\Sigma;\Delta \proves \fFunc{v_2} : {T_2}}$ implies
  $ \fType{\Sigma,\Sigma_2;\Delta,\Delta_2 \proves \fFunc{e_2[v_2/x_2]} : U}$,
  \end{itemize}
  then
  $ \fType{\Sigma_1,\Sigma_2;\Delta_1,\Delta_2 \proves
    \fFunc{\tCase{e}{x_1}{e_1}{x_2}{e_2}} : U}$. 
\end{lemma}

We exercise this extension with an example.

\begin{example}[Leader election]
The problem of leader election can be reduced to a process that makes
finite number of $\tau$ steps and then finishes with a $\fType{left}$
or $\fType{right}$ step on some channel of type $\fType{\fFunc{c} : \sTerm
  \ichoice \sTerm}$. To this end, we define the process $p_0 = (\set{\fFunc{cs},
  \fFunc{d}}, \mset{\fFunc{e_0}, \fFunc{e_1}, \fFunc{e_2}})$ which consists of
the following terms making use of the sum extension. We use standard syntactic
sugar for the encoding of the type $\fType{Bool} = \fType{\tUnit+\tUnit}$ and
the conditional $\fFunc{\mathit{if}}$.
%
\begin{align*}
  \fFunc{e_0} &=
                \fFunc{\tApp{\tClos}{(\tApp{\tSend}{\tPair{true}{\tApp{\tSend}{\tPair{false}{d}}}})}}
  \\
  \fFunc{e_1} &=
                \fFunc{\tLet[\tLVB][d]{\tApp{\tRecv}{d}}{\tLet*{\tApp{\tClos}{cs}}{\tIf{\tLVB}{\tApp{\tClos}{(\mathit{select}~\mathit{left}~c)}}{\tUnit}}}}
  \\
  \fFunc{e_2} &=
                \fFunc{\tLet[\tLVB][d]{\tApp{\tRecv}{d}}{\tLet*{\tApp{\tClos}{cs}}{\tLet*{\tApp{\tClos}{d}}{\tIf{\tLVB}{\tApp{\tClos}{(\mathit{select}~\mathit{right}~c)}}{\tUnit}}}}}
\end{align*}
Now we
set out to prove that $p_0 \in \lrProc{n}{ \fFunc{c}\fType{:  \sTerm
  \ichoice \sTerm} \vdash \fType{\tUnit}}$. This process has essentially
two different reduction sequences due to a race condition
(\rocq{leader_proc}). Either $\fFunc{e_0}$
communicates first with $\fFunc{e_1}$ and then with $\fFunc{e_2}$ or
vice versa. Each term, $\fFunc{e_1}$ and $\fFunc{e_2}$, receives a
boolean on channel $\fFunc{d}$, then synchronizes by closing
$\fFunc{cs}$, $\fFunc{e_2}$ closes $\fFunc{d}$, and we are left with
one term where the condition is $\fFunc{true}$ and another where it is
$\fFunc{false}$. The  $\fFunc{true}$-term executes a $\fFunc{select}$
and the $\fFunc{false}$-term does nothing.

The process $p_0$ is not reachable by a reduction sequence from a process of the
form $(\set{}, \mset{e})$, for some closed $e$, because
thread $\fFunc{e_0}$ shares $\fFunc{d}$ with both $\fFunc{e_1}$ and
$\fFunc{e_2}$. Term $\fFunc{e_1}$ is not typable
(cf. \cref{sec:syntacticTyping}) because it violates linearity for
variables/channels $\fFunc{d}$ and $\fFunc{c}$. Likewise, term
$\fFunc{e_2}$ violates linearity for channel $\fFunc{c}$.

See \cref{sec:reduct-lead-elect} in the appendix for the reduction sequence.
Composing $p_0$ over $\fFunc{c}$ with a syntactically typable client
yields a closed, deadlock-free process
(\rocq{leader_election_deadlock_free}).
\end{example}



\subsection{Recursion}
\label{sec:recursion}

\begin{figure}[tp]
  New values, terms, types \hfill \rocq{TBang}
  \begin{align*}
    \fFunc{v} &::= \dots \sepr \fFunc{\IExp{(\tIVX = e)}}
    & \text{recursively defined thunk}\\
    \fFunc{e} &::= \dots \sepr \fFunc{\tForce} & \text{thunk elimination}\\
    \fType{T} &::= \dots \sepr \fType{\tBox{T}}
  \end{align*}
  New labeled transition \hfill \rocq{HeadForce}
  \begin{mathpar}
    \begin{bussproof}[T-Act-Exp]
      \bussAx{\fFunc{\tApp{\tForce}{\IExp{(\tIVX = e)}}}
        \trans{\beta}
        \fFunc{\tSub{e}{\IExp{(\tIVX = e)}}{\tIVX}}}
    \end{bussproof}
  \end{mathpar}
  New typing rules for terms
  \hfill \rocq{typed_bang}
  \begin{mathpar}
    \begin{bussproof}[E-Exp]
      \bussAs{\fType{\tEmpty ; \tBox{\Delta} , \fFunc{\tIVX}:\tBox{T} \vdash \fFunc{e} : T}}
      \bussU{\fType{\tEmpty  ; \tBox{\Delta} \vdash \fFunc{\IExp{(\tIVX = e)}} : \tBox{T}}}
    \end{bussproof}
    \and
    \begin{bussproof}[E-Force]
      \bussAs{\fType{\tEmpty ; \tEmpty \vdash \fFunc{\tForce} : \lFunc{\tBox{T}}{T}}}
    \end{bussproof}
    \\
    \begin{bussproof}[E-Weak]
      \bussAs{\fType{\Sigma ; \Delta \vdash \fFunc{e} : U}}
      \bussU{ \fType{\Sigma ; \Delta, \tIVX : \tBox{T} \vdash \fFunc{e} : U}}
    \end{bussproof}
    \and
    \begin{bussproof}[E-Contract]
      \bussAs{\fType{\Sigma ; \Delta, \tIVX : \tBox{T}, \tIVY : \tBox{T} \vdash \fFunc{e} : U}}
      \bussU{ \fType{\Sigma ; \Delta, \tIVX : \tBox{T} \vdash \fFunc{\tSub{e}{\tIVX}{\tIVY}} : U}}
    \end{bussproof}
  \end{mathpar}
  \caption{Extending the term language and typing
    (\cref{fig:typeSys}) with recursive exponentials in the style of Lily}
  \Description{Additional values, terms and evaluation contexts for recursive exponentials; one reduction axiom; and three typing rules over a judgment carrying an extra intuitionistic environment.}
  \label{fig:lily:terms}
\end{figure}

We support term recursion by adapting ideas from Lily
\cite{DBLP:journals/entcs/BiermanPR00}, a linear lambda calculus with
recursion. Lily is based on a term calculus for linear logic with
exponentials where ``recursion is absorbed into terms of the
exponential type'' \cite{DBLP:journals/entcs/BiermanPR00}. In
consequence, we add a boxed type for duplicable values, an introduction form
for recursively defined values ``thunks'' in the style of Lily, and a $\fFunc{\tForce}$
operator as an eliminator for boxed types as shown in
\cref{fig:lily:terms}. Lily's eliminator is different, but equivalent.

The typing rules are as expected. Introducing a thunk requires an environment of
only boxed types, indicated by $\tBox{\Delta}$; forcing a thunk eliminates the
boxed type; boxed-typed variables are subject to weakening and contraction. The
mechanization implements the structural rules as part of environment splitting,
rather than explicit steps in a typing derivation.

Operationally, the thunk $\fFunc{\IExp{(\tIVX = e)}}$ is a recursively defined value and
the elimination operator exposes the body by unrolling the recursion.
% The reduction rule
% we propose here is proved to be a Kleene
% equivalence by \citet{DBLP:journals/entcs/BiermanPR00} with respect to
% their big-step semantics.
\citet{DBLP:journals/entcs/BiermanPR00} define some syntactic sugar for a fixed
point operator and non-recursive thunks:
\begin{align*}
  \fFunc{\tFixI{\tIVX}{e}} &:= \fFunc{\tApp{\tForce}{\IExp{(\tIVX = e)}}}
  & \fFunc{\IExp{e}} & := \fFunc{\IExp{(\tIVX = e)}} \quad \text{where}~\fFunc{\tIVX\notin fv (e)}
\end{align*}

We extend the value logical relation in a similar way as \citet{DBLP:journals/entcs/BiermanPR00}.
\begin{align*}
  \lrVal{n}{\fType{\tEmpty \vdash \tBox{T}}}
  &= \set{ \fFunc{\IExp{(\tIVX = e)}}  \mid
    \fc(\fFunc{e}) = \emptyset \land
    \fFunc{{\tFixI{\tIVX}{e}}} \in \lrTerm{n}{\fType{\tEmpty \vdash T}}
    }
\end{align*}
A replicable thunk owns no channels, which is what the side condition
$\fc(\fFunc{e}) = \emptyset$ records; exponentials are also admitted as
return types of a thread, so we add the rule
$\primitive \fType{T} \implies \primitive \fType{\tBox{T}}$
to the ones of \cref{sec:syntacticTyping:typeSys}.

For the additional cases of the fundamental theorem, we proved the
following compositionality lemmas, each corresponding to one of the typing rules
in \cref{fig:lily:terms}.

\begin{lemma}[Compositionality: \ruleLabel{exp} \rocq{sem_typed_bang}]
  \label{lem:comp:exp}
  If
  $ \fType{\tEmpty;\tBox{\Delta}, \tBox{T} \proves \fFunc{e} : T}$,
  then
  $ \fType{\tEmpty;\tBox{\Delta} \proves \fFunc{\IExp{(x=e)}} : T}$.
\end{lemma}

\begin{lemma}[Compositionality: \ruleLabel{force} \rocq{sem_typed_force}]
  \label{lem:comp:force}
  If
  $ \fType{\Sigma; \Delta \proves \fFunc{e} : \tBox{T}}$,
  then
  $ \fType{\Sigma; \Delta \proves \fFunc{\tApp{\tForce}{e}} : {T}}$.
\end{lemma}

\begin{lemma}[Compositionality: \ruleLabel{weak} \rocq{sem_typed_wk}]
  If
  $ \fType{\Sigma; \Delta \proves \fFunc{e} : U}$,
  then
  $ \fType{\Sigma; \Delta, x:\tBox{T} \proves \fFunc{e} : U}$.
\end{lemma}
\begin{lemma}[Compositionality: \ruleLabel{contract} % \rocq{sem_typed_contract}
  ]
  If
  $ \fType{\Sigma; \Delta, x:\tBox{T}, y:\tBox{T} \proves \fFunc{e} : U}$,\\
  then
  $ \fType{\Sigma; \Delta, x:\tBox{T} \proves \fFunc{\tSub{e}{x}{y}} : U}$.
\end{lemma}

This encoding of recursion also supports replication as in CP
\cite{conf/icfp/Wadler12,DBLP:journals/jfp/Wadler14}, the process calculus
companion to GV \cite{conf/esop/LindleyM15}. We implement a replicable
unrestricted session carrying type $\fType{S}$ as a function
of type $\fType{[S]} := \fType{\tBox{(\lFunc{\tUnit}{S})}}$ and
treat the new operators as macros as follows.
\begin{align*}
  \fFunc{repl} & \fType{{}:{\lFunc{\tBox{(\lFunc{S}{\tUnit})}}{[{\dual{S}}]}}}
  &
    \fFunc{req} & \fType{{}:{\lFunc{[S]}{S}}} \\
  \fFunc{repl} & :=
                 \fFunc{\tAbs{x}{\IExp{(\tAbs{\_}{\tApp{\tFork}{(\tApp{\tForce}{x})}})}}}
  & \fFunc{req} & := \fFunc{\tAbs{x}{\tApp{(\tApp{\tForce}{x})}{\tUnit}}}
\end{align*}

\mw{claude: likewise here --- examples/Factory.v mechanizes a
  replicated server: one factory value serving two independent
  clients, plus a discarded request, typed and deadlock free. It is
  the evidence that the encoding of Sect. 4.4 is usable end to end.
  Add it as an example?}
\pt{Yes, add this example here}

\subsection{Recursive Types}
\label{sec:extens-recurs-ii}

In the presence of recursive functions, it makes sense to extend
session types to be recursive, too. Specifically, we adopt the traditional approach of
equi-recursive session types, where a recursive type is always equal to its unfolding.
\begin{align*}
  \fType{S} &::= \dots \sepr \fType{\mu\alpha.S} \sepr \fType{\alpha}
  & \fType{\mu\alpha.S} &= \fType{\tSub{S}{\mu\alpha.S}{\alpha}}
\end{align*}
As equi-recursive types have contractive definitions
(\rocq{contractive_ty}), we can always
determine the top-level constructor of a type in finite time
(\rocq{unfoldS}). Hence,
all our prior definitions and proofs carry over without change. In the
mechanization, we insert appropriate unfolding functions that expose the
top-level constructor for inspection of the respective clause in the logical
relation.

\mw{claude: the mechanization has a worked example for this section
  that the paper does not: a recursive ping/pong over an equi-recursive
  session type, typed and deadlock free (examples/Pingpong.v). Add it
  as a short example, the way Sect. 4.1 has the leader election?}
\pt{this comment was for the previous subsection on recursion, but it makes more
sense to add it here where we have recursive types!}

\subsection{The Link Operator}
\label{sec:extens-link-oper}

\begin{figure}[tp]
  New constant, session label \hfill \rocq{ELink}
  \begin{align*}
    \fFunc{k} &::= \dots \sepr \fFunc{\tLink}
    & \text{forwarding operator}
    \\
    \sigma &::= \dots \sepr \lLink{c}{c}
             &\text{forwarding session label}
  \end{align*}
  New labeled transitions \hfill \rocq{HeadLink}
  \begin{mathpar}
    \begin{bussproof}[T-Act-Link]
      \bussAx{\fFunc{\tApp{\tLink}{\tPair{c}{d}}} \trans{\lLink{c}{d}} \fFunc{\tUnit}}
    \end{bussproof}
        \and
    \begin{bussproof}[P-Act-Link]
      \bussAs{\fFunc{e} \trans{\lLink{c}{d}} \fFunc{e'}}
      \bussU{\fProc{(C \cupdot \set{\fFunc{c}}, E \uplus \mset{\fFunc{e}})} \trans{\tau} \fProc{(C, \tSub{(E \uplus \mset{\fFunc{e'}})}{\fFunc{d}}{\fFunc{c}})}}
    \end{bussproof}
  \end{mathpar}
  Subjects and escaping channels of the new label:
  $ \fc(\lLink{c}{d}) = \set{\fFunc{c},\fFunc{d}} $, \quad
  $ \sent(\lLink{c}{d}) = \emptyset $

  New type-indexed family of constants: \hfill \rocq{KTLink}\\
  $ \fType{\fFunc{\tLink} : \lFunc{\tPair{S}{\dual{S}}}{\tUnit}} $

  New case for the term logical relation \hfill \rocq{Edef}
  \begin{enumerate}[start=6]
  \item  $\lambda = \lLink{c}{d}$ implies $\fFunc{c} \ne \fFunc{d}$ and $\fType{\Sigma} =
    \fType{\Sigma_1, \fFunc{c}:S, \fFunc{d}:\dual{S}}$ and $\fFunc{e'} \in \lrTerm{n'}{\fType{\Sigma_1 \vdash T}}$.
  \end{enumerate}

  New case for the process logical relation \hfill \rocq{Pdef}
  \begin{enumerate}[start=5]
  \item  $\pi = \lLink{c}{d}$ implies $\fFunc{c} \ne \fFunc{d}$ and $\fType{\Sigma} =
    \fType{\Sigma_1, \fFunc{c}:S, \fFunc{d}:\dual{S}}$ and $\fProc{p'} \in \lrProc{n'}{\fType{\Sigma_1 \vdash T}}$.
  \end{enumerate}

  The same two cases as obligation rules
  (\cref{fig:logrel-rules-term,fig:logrel-rules-proc})
  \begin{mathpar}
    \begin{bussproof}[O-Link]
      \bussAs{\fFunc{c} \ne \fFunc{d}}
      \bussAs{\fFunc{e'} \in \lrTerm{n}{\fType{\Sigma_1 \vdash T}}}
      \bussB{\obl{\fType{\Sigma_1, \fFunc{c}:S, \fFunc{d}:\dual{S}}}{n}{\fFunc{e} \trans{\lLink{c}{d}} \fFunc{e'} : \fType{T}}}
    \end{bussproof}
    \and
    \begin{bussproof}[Q-Link]
      \bussAs{\fFunc{c} \ne \fFunc{d}}
      \bussAs{\fProc{p'} \in \lrProc{n}{\fType{\Sigma_1 \vdash T}}}
      \bussB{\obl{\fType{\Sigma_1, \fFunc{c}:S, \fFunc{d}:\dual{S}}}{n}{\fProc{p} \trans{\lLink{c}{d}} \fProc{p'} : \fType{T}}}
    \end{bussproof}
  \end{mathpar}

  \caption{Extensions for forwarding (the link operator)}
  \Description{The new link constant and forwarding session label; the term and process reduction rules for forwarding; the subjects and escaping channels of the new label; the type of link; and the new case of the term and process logical relations, given both as a clause and as an inference rule.}
  \label{fig:link-operator}
\end{figure}

Session type systems based on the Curry-Howard correspondence for
classical linear logic (e.g.,
\cite{DBLP:journals/jfp/Wadler14,conf/esop/LindleyM15}) often include
a forwarding operator. This operator (called $\tLink$ by
\citet{conf/esop/LindleyM15}) takes two channel ends with dual
session types and connects them. This way, all messages on one channel
are forwarded to the other channel, in both directions.

\Cref{fig:link-operator} shows how to extend $\OurLang$ with
forwarding. We add a new operator $\fFunc{\tLink}$ and a forwarding label
$\lLink{\fFunc{c}}{\fFunc{d}}$ to indicate the request to forward between sessions
$\fFunc{c}$ and $\fFunc{d}$. As forwarding is symmetric the label
carries a \emph{set} of two channels, i.e., their order does not matter.
The two ends are \emph{not} required to be different: a self-forwarding
$\fFunc{\tApp{\tLink}{\tPair{c}{c}}}$ is a redex like any other and
emits the singleton label $\lLink{c}{c} = \set{\fFunc{c}}$.
The label is generated by a term
reduction. The corresponding process reduction requires that one of
the channel ends involved in forwarding is local. In this case, one
local channel ($\fFunc{c}$) is removed and the other channel is
substituted for $\fFunc{c}$. If no local channel is involved in the
label, the configuration as a whole takes a labeled reduction step
(cf. \TirName{P-Act-Session} in \cref{fig:proc-semantics}).
The new cases for the logical relations formalize the postconditition
of the forwarder: the two connected channel ends vanish from the
context as a consequence of the forwarding reduction.
Distinctness of the two ends is imposed by the logical relation, and
only there: a configuration offering a forwarding whose two ends
coincide names a single channel, which no instance of the case above
matches, so such a configuration is simply not related.
No typed term builds one either, since $\fFunc{\tLink}$ consumes a
session type together with its dual and linearity forbids one channel
supplying both.
On a restricted channel the substitution $\tSub{\cdot}{\fFunc{c}}{\fFunc{c}}$
is the identity, so \ruleLabel{P-Act-Link} degenerates to releasing the
name --- which is what a self-forwarding means: the thread held both
ends, so the session is spent.

\mw{claude: the mechanization carries a worked example for this
  section that the paper does not have --- Milner's cyclic scheduler
  with one agent replaced by a forwarder (Scheduler.v), closed,
  deadlock free, and untypable in Curry-Howard systems (it is
  Example~1 of Dardha and Gay). Add it as a subsection?}
Installing this extension, requires to update the proof of
\ruleLabel{E-Abs}
to cater for the new functional constant and
all other rules to cater for the new labeled reduction.

%%% Local Variables:
%%% mode: latex
%%% TeX-master: "main"
%%% End:



\section{Mechanization}
\label{sec:mechanization}

\mw{claude: emptied. What stood here described the Agda prototype ---
  intrinsically scope-safe terms over $Fin~n$, Autosubst2-style
  substitutions, soups as lists with permutation invariance, resource
  vectors instead of context splitting --- and none of it describes the
  current artifact, which is a Rocq development on top of stdpp
  (\texttt{gset}/\texttt{gmultiset}, so soups really are sets and
  multisets, and contexts really are split). It is axiom free and
  postulates nothing. Rewrite from scratch, or drop the section and
  fold a paragraph into the introduction?}

%%% Local Variables:
%%% mode: latex
%%% TeX-master: "main"
%%% End:


\section{Related Work and Future Directions}
\label{sec:relatedWork}


In this paper, we presented the first logical relation for session types
which supports GV-style rules, can prove deadlock freedom, and
admits safe processes which do not comply with communication topology
constraints imposed by the syntactic type system.
In this section we review the literature in view of the technical development
and sketch potential future developments.

\citet{DBLP:journals/iandc/PerezCPT14} were the first to investigate
session types using logical relations.
Their starting point is the logical account of session types \cite{DBLP:conf/concur/CairesP10}; their process calculus includes a (non-recursive)
treatment of the exponential, and uses the logical relation to prove
strong normalization, confluence, and behavioral equality.
More recent work by \citet{conf/esop/RochaC23} uses a logical relation that specifies
strong normalization and safe handling of state for message-passing
processes using linear session types
(based on an extension of ILL~\cite{conf/tapsoft/GirardL87}).
In both \cite{DBLP:journals/iandc/PerezCPT14,conf/esop/RochaC23},
the logical relation is defined only for syntactically typed processes.
Deadlock freedom is established by proving subject reduction and progress
on the syntactic derivations, and the proofs of the logical relation
rely on having externally established deadlock freedom.
In contrast, we establish all the behavioral guarantees of the type system
as a consequence of the semantic logical relation.

An interesting direction for further work would be to establish
compositional and liberal logical relations for properties beyond
deadlock freedom.
For example,
\citet{conf/esop/Atkey17} defines an operational semantics for
Wadler's ``processes as CLL proofs''~\cite{conf/icfp/Wadler12} and
uses a logical relation to prove adequacy
of the operational semantics and a type-based
denotational semantics with the goal
of establishing observational equivalences.
\citet{DBLP:conf/lics/DerakhshanBJ21} study noninterference for
information flow control.
Van den Heuvel et al.~\cite{DBLP:conf/ecoop/0001DB24} extend the
work to cyclic processes.
\citet{DBLP:journals/pacmpl/YaoICBZ25} consider timed
message passing for IoT systems via session types defined by a logical relation.

The logical-relations approach also offers a nice opportunity
when extending the language of types to support polymorphism,
as it would allow the derivation of parametricity theorems from polymorphic types.


Another direction would be to extend the features of the programming model.
For example, LinearActris~\cite{JacobsHK24} supports mutable heaps
in addition to message passing primitives.
It would be interesting to try and embed our logical relation into
a separation logic, allowing a more flexible mix of paradigms.

While this work focused on the Curry-Howard based ways of composing processes,
approaches based on priorities, \eg \cite{Kobayashi97,Padovani14,DardhaG18},
track dependencies between actions across sessions explicitly in the types.
Finding a semantic logical relation capturing priority-based reasoning
would provide powerful compositional means of proving well-behavedness of
components that cannot be syntactically typed (like our load balancer example).
It could also provide a semantic bridge between priority-based and Curry-Howard-based systems, in the spirit of \citet{DardhaP22}.

Multiparty Session Types also shed light on global dependencies in types.
For programs interleaving several multiparty sessions,
establishing deadlock freedom requires additional priority-based annotations and checks.
If one considers programs with a single session involving multiple participants,
the Multiparty Session Type discipline ensures deadlock freedom.
The semantic logical relation understanding of Multiparty Session Types is
still underdeveloped.
The Multris logic~\cite{HinrichsenJK24} made some steps in developing
a logical relation embedded into the Iris separation logic,
but it can only prove safe usage of channels and resources,
not deadlock freedom.

Finally, a challenging open problem is to provide a semantic
and compositional logical relation to prove \emph{liveness} properties
like lock freedom.
Priority-based systems have been adapted to prove lock freedom and progress
properties~\cite{Kobayashi02,Padovani13},
but the soundness argument is syntactic.
Several works in Concurrent Separation Logics 
\cite{DOsualdoSFG21,LiangF18,LeeLYCKH25}
explored compositional semantic characterizations of progress properties
for shared memory concurrent programs.
Combining those ideas with our logical relation approach
could provide some leverage to attack the problem.

%%% Local Variables:
%%% mode: latex
%%% TeX-master: "main"
%%% End:



\bibliographystyle{splncs04nat} % https://github.com/tpavlic/splncs04nat
\bibliography{main.bib}

\clearpage
\appendix

% \section{Safe process with payload}
% \label{sec:safe-process-with}

% The following process is safe and in $\lrProc{n}{\fType{\tEmpty} \vdash \fType{\tUnit}}$.

% \begin{align*}
%   & (\set{},\mset{ \fFunc{
%     \tLet[\tLVB][\tLVA]{\tApp{\tRecv}{(\tApp{\tFork}{\tAbs{\tLVA}{\tApp{\tClos}{(\tApp{\tSend}{\tPair{\tUnit}{\tLVA}})}}})}}{\tLet*{\tLVB}{\tApp{\tClos}{\tLVA}}}
%     } })
%   \\
%   \redu
%   & (\set{ \fFunc{c} },\mset{ \fFunc{
%     \tLet[\tLVB][\tLVA]{(\tApp{\tRecv}{c})}{\tLet*{\tLVB}{\tApp{\tClos}{\tLVA}}}
%     }, \fFunc{
%     \tApp{(\tAbs{\tLVA}{\tApp{\tClos}{(\tApp{\tSend}{\tPair{\tUnit}{\tLVA}})}})}{c}
%     } })
%   \\
%   \redu
%   & (\set{ \fFunc{c} },\mset{ \fFunc{
%     \tLet[\tLVB][\tLVA]{(\tApp{\tRecv}{c})}{\tLet*{\tLVB}{\tApp{\tClos}{\tLVA}}}
%     }, \fFunc{
%     {\tApp{\tClos}{(\tApp{\tSend}{\tPair{\tUnit}{c}})}}
%     } })
%   \\
%   \redu
%   & (\set{ \fFunc{c} },\mset{ \fFunc{
%     \tLet[\tLVB][\tLVA]{\tPair{\tUnit}{c}}{\tLet*{\tLVB}{\tApp{\tClos}{\tLVA}}}
%     }, \fFunc{
%     {\tApp{\tClos}{c}}
%     } })
%   \\
%   \redu
%   & (\set{ \fFunc{c} },\mset{ \fFunc{
%     {\tLet*{\tUnit}{\tApp{\tClos}{c}}}
%     }, \fFunc{
%     {\tApp{\tClos}{c}}
%     } })
%   \\
%   \redu
%   & (\set{ \fFunc{c} },\mset{ \fFunc{
%     {{\tApp{\tClos}{c}}}
%     }, \fFunc{
%     {\tApp{\tClos}{c}}
%     } })
%   \\
%   \redu
%   & (\set{ },\mset{ \fFunc{
%     \tUnit
%     }, \fFunc{
%     \tUnit
%     } })
% \end{align*}

% We have $\fFunc{
%   \tLet[\tLVB][\tLVA]{\tApp{\tRecv}{(\tApp{\tFork}{\tAbs{\tLVA}{\tApp{\tClos}{(\tApp{\tSend}{\tPair{\tUnit}{\tLVA}})}}})}}{\tLet*{\tLVB}{\tApp{\tClos}{\tLVA}}}
% } \in \lrTerm{n}{\fType{\tEmpty} \vdash \fType{\tUnit}}$ because
% $\tLet[\tLVB][\tLVA]{(\tApp{\tRecv}{c})}{\tLet*{\tLVB}{\tApp{\tClos}{\tLVA}}}
% \in \lrTerm{n}{\fFunc{c}\fType{: \sRecv{\tUnit}{\tClos}} \vdash
%   \fType{\tUnit}}$.
% The reduction sequence is similar to the one shown for the first
% example in \cref{sec:safe-process}.

\section{Reductions for leader election example}
\label{sec:reduct-lead-elect}

As an example, we trace the reduction sequence starting from $p_0$  where $\fFunc{e_0}$
communicates first with $\fFunc{e_1}$
(\rocq{leader_proc}, \rocq{leader_sem_ptyped}).
\begin{itemize}
\item $\tau$ transition (internal communication) to
  \begin{small}
    \begin{align*}
      \fFunc{e_0} &=
                    \fFunc{\tApp{\tClos}{(\tApp{\tSend}{\tPair{true}{d}})}}
      \\
      \fFunc{e_1} &=
                    \fFunc{\tLet[\tLVB][d]{\tPair{false}{d}}{\tLet*{\tApp{\tClos}{cs}}{\tIf{\tLVB}{\tApp{\tClos}{(select~left~c)}}{\tUnit}}}}
      \\
      \fFunc{e_2} &=
                    \fFunc{\tLet[\tLVB][d]{\tApp{\tRecv}{d}}{\tLet*{\tApp{\tClos}{cs}}{\tLet*{\tApp{\tClos}{d}}{\tIf{\tLVB}{\tApp{\tClos}{(select~right~c)}}{\tUnit}}}}}
    \end{align*}
  \end{small}
\item $\tau$ transition ($\beta$ pair elimination) to
  \begin{small}
    \begin{align*}
      \fFunc{e_0} &=
                    \fFunc{\tApp{\tClos}{(\tApp{\tSend}{\tPair{true}{d}})}}
      \\
      \fFunc{e_1} &=
                    \fFunc{{\tLet*{\tApp{\tClos}{cs}}{\tIf{false}{\tApp{\tClos}{(select~left~c)}}{\tUnit}}}}
      \\
      \fFunc{e_2} &=
                    \fFunc{\tLet[\tLVB][d]{\tApp{\tRecv}{d}}{\tLet*{\tApp{\tClos}{cs}}{\tLet*{\tApp{\tClos}{d}}{\tIf{\tLVB}{\tApp{\tClos}{(select~right~c)}}{\tUnit}}}}}
    \end{align*}
  \end{small}
\item $\tau$ transition (internal communication) to
  \begin{small}
    \begin{align*}
      \fFunc{e_0} &=
                    \fFunc{\tApp{\tClos}{d}}
      \\
      \fFunc{e_1} &=
                    \fFunc{{\tLet*{\tApp{\tClos}{cs}}{\tIf{false}{\tApp{\tClos}{(select~left~c)}}{\tUnit}}}}
      \\
      \fFunc{e_2} &=
                    \fFunc{\tLet[\tLVB][d]{\tPair{true}{d}}{\tLet*{\tApp{\tClos}{cs}}{\tLet*{\tApp{\tClos}{d}}{\tIf{\tLVB}{\tApp{\tClos}{(select~right~c)}}{\tUnit}}}}}
    \end{align*}
  \end{small}
\item $\tau$ transition ($\beta$ pair elimination) to
  \begin{small}
    \begin{align*}
      \fFunc{e_0} &=
                    \fFunc{\tApp{\tClos}{d}}
      \\
      \fFunc{e_1} &=
                    \fFunc{{\tLet*{\tApp{\tClos}{cs}}{\tIf{false}{\tApp{\tClos}{(select~left~c)}}{\tUnit}}}}
      \\
      \fFunc{e_2} &=
                    \fFunc{{\tLet*{\tApp{\tClos}{cs}}{\tLet*{\tApp{\tClos}{d}}{\tIf{true}{\tApp{\tClos}{(select~right~c)}}{\tUnit}}}}}
    \end{align*}
  \end{small}
\item $\tau$ transition (internal communication, closing $\fFunc{cs}$,
  new $C = \set{\fFunc{c}, \fFunc{d}}$)
  \begin{small}
    \begin{align*}
      \fFunc{e_0} &=
                    \fFunc{\tApp{\tClos}{d}}
      \\
      \fFunc{e_1} &=
                    \fFunc{{\tLet*{\tUnit}{\tIf{false}{\tApp{\tClos}{(select~left~c)}}{\tUnit}}}}
      \\
      \fFunc{e_2} &=
                    \fFunc{{\tLet*{\tUnit}{\tLet*{\tApp{\tClos}{d}}{\tIf{true}{\tApp{\tClos}{(select~right~c)}}{\tUnit}}}}}
    \end{align*}
  \end{small}
\item $\tau$ transitions (2x $\beta$ unit)
  \begin{small}
    \begin{align*}
      \fFunc{e_0} &=
                    \fFunc{\tApp{\tClos}{d}}
      \\
      \fFunc{e_1} &=
                    \fFunc{{{\tIf{false}{\tApp{\tClos}{(select~left~c)}}{\tUnit}}}}
      \\
      \fFunc{e_2} &=
                    \fFunc{{{\tLet*{\tApp{\tClos}{d}}{\tIf{true}{\tApp{\tClos}{(select~right~c)}}{\tUnit}}}}}
    \end{align*}
  \end{small}
\item $\tau$ transition (internal communication, closing $\fFunc{d}$,
  new $C = \set{\fFunc{c}}$)
  \begin{small}
    \begin{align*}
      \fFunc{e_0} &=
                    \fFunc{\tUnit}
      \\
      \fFunc{e_1} &=
                    \fFunc{{{\tIf{false}{\tApp{\tClos}{(select~left~c)}}{\tUnit}}}}
      \\
      \fFunc{e_2} &=
                    \fFunc{{{\tLet*{\tUnit}{\tIf{true}{\tApp{\tClos}{(select~right~c)}}{\tUnit}}}}}
    \end{align*}
  \end{small}
\item $\tau$ transition (unit elim)
  \begin{small}
    \begin{align*}
      \fFunc{e_0} &=
                    \fFunc{\tUnit}
      \\
      \fFunc{e_1} &=
                    \fFunc{{{\tIf{false}{\tApp{\tClos}{(select~left~c)}}{\tUnit}}}}
      \\
      \fFunc{e_2} &=
                    \fFunc{{{{\tIf{true}{\tApp{\tClos}{(select~right~c)}}{\tUnit}}}}}
    \end{align*}
  \end{small}
\item $\tau$ transitions (2x $\beta$ bool)
  \begin{small}
    \begin{align*}
      \fFunc{e_0} &=
                    \fFunc{\tUnit}
      \\
      \fFunc{e_1} &=
                    \fFunc{{{{\tUnit}}}}
      \\
      \fFunc{e_2} &=
                    \fFunc{{{{{\tApp{\tClos}{(select~right~c)}}}}}}
    \end{align*}
  \end{small}
\item There are two remaining labeled transitions: $\lSend{c}{right}$
  followed by $\lTerm{c}$. The resulting process is terminated and is
  an element of $\lrProc{n}{\fType{\tEmpty} \vdash \fType{\tUnit}}$:
  \begin{small}
    \begin{align*}
      \fFunc{e_0} &=
                    \fFunc{\tUnit}
      \\
      \fFunc{e_1} &=
                    \fFunc{{{{\tUnit}}}}
      \\
      \fFunc{e_2} &=
                    \fFunc{{{{{\tUnit}}}}}
    \end{align*}
  \end{small}
\end{itemize}

\end{document}
